\section{The a-posteriori theorem at fixed frequency (G1)}\label{bp:p2}

This part proves Theorem~\ref{bp:p0:g1}: an approximately invariant torus with nondegenerate frame and
twist, whose invariance error is small compared with $\gamma^4$, lies within $O(e/\gamma^2)$ of a true
invariant torus with the same Diophantine frequency. The proof is the quasi-Newton scheme of HL19
(Theorem~3.1, pp.~16--18; \S4--5; Appendix~B), specialized as follows.
\begin{itemize}
\item $d=n$: there are no first integrals, the tangent frame is $L=DK$, and every $X_p$-term of HL19
  vanishes (HL19 Rem.~3.2, p.~18).
\item The canonical structure $\Omega=\tilde\Omega=\Omega_0$, $G=I_{2n}$, $J=\Omega_0$ (HL19 Rem.~3.3,
  p.~18), in Case~III, where $A=0$ (HL19 \S2.5, p.~12; (3.5), p.~16). Then $B=(L^\top L)^{-1}$ and
  $N=\Omega_0LB$, which are the frozen definitions of Part~\ref{bp:p0}.
\item Primary tori in $\T^n\times U$ (HL19 Rem.~2.3, p.~7): $K(\theta+k)=K(\theta)+(k,0)$, so $DK$, the
  frame, the error and the torsion are $1$-periodic. Lemma~\ref{bp:p2:periodicity} checks every formula
  we use against this.
\item Sup norms over \emph{open} strips, stated pointwise (decision D12): ``$\abs f\le F$ on $T_r$''
  means $\abs{f(\theta)}\le F$ for every $\theta\in T_r$. No closed strip and no continuity to the
  boundary is used, so no extra shrinking step is needed.
\item The cohomological equation is solved with the crude exponent $\varkappa=\tau+n$ of
  Lemma~\ref{bp:p1:cohomological} instead of R\"ussmann's $\tau$. Only powers of the strip losses
  $\delta$ change; the scheme stays quadratic and the powers of $\gamma$ are those of HL19.
\end{itemize}
GJLV05 (Theorem~5, p.~865; \S8, pp.~889--893) is the cross-check for flows.

\paragraph{Strategy.} The tangent vectors $L=DK$ and the normal vectors $N=\Omega_0LB$ form a frame
$P=(L\ N)$ along the torus. If $K$ were invariant, $P$ would be symplectic and would conjugate the
linearized vector field $DX_h(K)-\partial_\omega$ to the triangular matrix
$\Lambda=\left(\begin{smallmatrix}0&T\\0&0\end{smallmatrix}\right)$, $T$ the torsion. For an
approximately invariant torus both properties hold up to errors controlled by the invariance error~$E$:
the symplecticity defect is the pull-back $\Omega_K=L^\top\Omega_0L$ of the symplectic form, which is
$O(e/(\gamma\delta^{\varkappa+1}))$ because it is a zero-average solution of a cohomological equation
with right side $O(e/\delta)$ (Subsection~\ref{bp:p2:sec-geometry}); the reducibility defect is
$O(e/\delta)$, with no small divisor at all in the canonical Case~III. The Newton correction
$\Delta K=P\xi$ is found by solving the triangular system, which needs two cohomological equations and
one inversion of the averaged torsion $\avg T$ (Subsection~\ref{bp:p2:sec-correction}). The new error
is quadratic, $O(e^2/(\gamma^4\delta^{2a}))$ with $a=2\varkappa+2$, and the frame, the Gram inverse and
the twist change by $O(e/(\gamma^2\delta^{a+2}))$ (Subsections~\ref{bp:p2:sec-perturbation}
and~\ref{bp:p2:sec-step}). With strip losses $\delta_s=2^{-s}\delta_0$ the errors decay
super-exponentially and all changes are summable; the margins~$\mu$ of the hypotheses absorb the sums
and leave the margins $\mu/2$ of the conclusion (Subsections~\ref{bp:p2:sec-iteration}
and~\ref{bp:p2:sec-proof}).

\paragraph{What differs from HL19.} Besides the specializations above: (i) the frequency is bounded by
the hypotheses, $\abs\omega\le2n\sigma_K\sigma_B(M+e)$ (Lemma~\ref{bp:p2:frequency}); this replaces
HL19's operator calculus for $L_\omega$ in Lemma~4.8 and Step~6 of Lemma~5.2 by plain Cauchy estimates,
at the cost of powers of $\delta$ only; (ii) the Diophantine condition at $k=e_1$ gives
$\gamma\le\abs{\omega}$, so the step constants do not depend on $\gamma$ (in HL19 they increase with
$\gamma$, ledger D-G1-1); the bound $\bar\gamma$ is used once, to force $e\le1$; (iii) constants are
crude polynomials in one size constant~$\mathfrak m$. All deviations are listed in
Subsection~\ref{bp:p2:sec-ledger}.

%----------------------------------------------------------------------------------------------------
\subsection{Standing assumptions, norms and constants}\label{bp:p2:sec-setting}

\begin{assumption}[Standing data of Part~\ref{bp:p2}]\label{bp:p2:standing}
Fix $n\ge2$, $\tau\ge n-1$, and $M,\sigma_K,\sigma_B,\sigma_T>0$. Let
$\mathfrak B\subset\C^n\times\C^n$ be open and invariant under $(\theta,I)\mapsto(\theta+k,I)$ for every
$k\in\Z^n$. Let $h$ be holomorphic on $\mathfrak B$, $1$-periodic in $\theta$
($h(\theta+k,I)=h(\theta,I)$ on $\mathfrak B$), real at real points of $\mathfrak B$, and
\[
\abs{X_h(z)}\le M,\qquad\abs{DX_h(z)}\le M,\qquad\abs{D^2X_h(z)}\le M\qquad(z\in\mathfrak B),
\]
where $D^2X_h(z)$ carries the norm of bilinear maps. Let $\gamma>0$ and $\omega\in\Dio_{\gamma,\tau}$.
\end{assumption}

The hypotheses of Theorem~\ref{bp:p0:g1} contain these, together with $\abs{Dh}\le M$, which is not
used. Only $\tau\ge0$ and $n\ge1$ are used below; $n\ge2$ and $\tau\ge n-1$ are inherited from the
frozen statement.

\paragraph{Norms.} As in Part~\ref{bp:p0}: vectors carry the max norm, with
$\abs{(u,v)}=\max\{\abs u,\abs v\}$ on $\C^n\times\C^n$; a $p\times q$ matrix $M$ carries the operator
norm $\abs M=\max_i\sum_j\abs{M_{ij}}$; bilinear maps carry the operator norm. The strip $T_r$ is open.
We write $\partial_\omega f:=Df\,\omega=\sum_\ell\omega_\ell\,\partial_{\theta_\ell}f$ for the derivative
of a function of $\theta$ along~$\omega$, and $\mathbb B_\beta(z)=\{w\in\C^n\times\C^n:\abs{w-z}<\beta\}$
for open balls. The identity matrix of size $m$ is $I_m$.

\begin{definition}[Exponents and constants]\label{bp:p2:constants}
Let $c(n,\tau)$ be the constant of Lemma~\ref{bp:p1:cohomological}. Define
\begin{gather*}
\varkappa:=\tau+n,\qquad a:=2\varkappa+2,\qquad c_R:=\max\{1,c(n,\tau)\},\\
\mathfrak m:=\max\{1,M,\sigma_K,\sigma_B,\sigma_T\},\qquad \Theta:=4n\mathfrak m^3,
\end{gather*}
and, in this order,
\begin{align*}
C_X&:=5n^2\mathfrak m^5, & C_T&:=2n\mathfrak m^2C_X,\\
C_\Omega&:=2n(n+1)c_R\mathfrak m, & C_{\mathrm{sym}}&:=\mathfrak m^2C_\Omega,\\
C_{\mathrm{red}}&:=2n(n+1)\mathfrak m^2, & C_{\xi_0}&:=2n\mathfrak m^2(\mathfrak m\Theta+c_RC_T),\\
C_{\xi N}&:=C_{\xi_0}+2nc_R\mathfrak m, & C_g&:=C_TC_{\xi N}+2n\mathfrak m^2\Theta,\\
C_\xi&:=c_RC_g+\Theta C_{\xi N}, & C_{\Delta K}&:=2\mathfrak m^2C_\xi,\\
C_{\mathrm{lin}}&:=C_{\mathrm{red}}C_\xi+C_{\mathrm{sym}}C_g, &
  C_E&:=4\mathfrak m^2\Theta^2C_{\mathrm{lin}}+\mathfrak mC_{\Delta K}^2,\\
C_G&:=6n^2\mathfrak m, & C_B&:=2\mathfrak m^2C_G,\\
C_{\Delta N}&:=2n\mathfrak m+\mathfrak mC_B, & C_{\Delta X}&:=2\mathfrak m^3+(\mathfrak m+n\Theta)C_{\Delta N},\\
C_{\Delta T}&:=4n\mathfrak m^2C_{\Delta X}+2nC_{\Delta N}C_X, &
  \mathfrak c_{\mathrm{step}}&:=2\Theta^4\max\{C_{\mathrm{sym}},\,2\mathfrak mC_{\Delta T}C_{\Delta K}\},\\
\mathfrak c_{\mathrm{sum}}&:=4\Theta^2C_{\Delta K}\max\{n,\,C_B,\,\mathfrak m^2C_{\Delta T}\}.
\end{align*}
All of these depend only on $(n,\tau,M,\sigma_K,\sigma_B,\sigma_T)$ and are $\ge1$. For $\mu>0$ and
$0<\rho_\infty<\rho$ define further
\[
\mathfrak c_*:=\max\Bigl\{2^{2a}C_E,\ \mathfrak c_{\mathrm{step}}\max\{1,2/\mu\},\
\frac{2\mathfrak c_{\mathrm{sum}}}{\mu}\Bigr\},\qquad
\delta_0:=\min\Bigl\{1,\frac{\rho-\rho_\infty}{8}\Bigr\},
\]
so that $\mathfrak c_*=\mathfrak c_*(n,\tau,M,\sigma_K,\sigma_B,\sigma_T,\mu)$ and
$\delta_0=\delta_0(\rho,\rho_\infty)$.
\end{definition}

The numerical factors are crude; only the dependence matters. Since $\tau\ge0$ and $n\ge1$ we have
$\varkappa\ge1$ and $a\ge4$; with the frozen $n\ge2$, $\tau\ge1$ we even have $a\ge8$. We use $a\ge3$
in Lemma~\ref{bp:p2:error-decay}.

%----------------------------------------------------------------------------------------------------
\subsection{Tools}\label{bp:p2:sec-tools}

This subsection collects elementary facts used throughout: matrix norms, Cauchy estimates along
$\omega$, boundedness and averages of periodic functions, the cohomological equation in operator norms,
and the Hamiltonian identities. Results not in Part~\ref{bp:p1} are listed for transfer in
Subsection~\ref{bp:p2:sec-transfer}.

\begin{lemma}[Matrix norms]\label{bp:p2:matrix}
\begin{enumerate}[label=(\alph*)]
\item For a $p\times q$ matrix $M$: $\max_{i,j}\abs{M_{ij}}\le\abs M=\sup_{\abs v\le1}\abs{Mv}$, and
  $\abs{M^\top}\le p\abs M$ if $q\le p$ (all transposes below are of $2n\times n$ matrices).
\item $\abs{M_1M_2}\le\abs{M_1}\abs{M_2}$ and $\abs{Mv}\le\abs M\abs v$.
\item For block matrices, $\abs{(M_1\ M_2)}\le\abs{M_1}+\abs{M_2}$,
  \[
  \Bigl|\begin{pmatrix}M_{11}&M_{12}\\M_{21}&M_{22}\end{pmatrix}\Bigr|
  \le\max\{\abs{M_{11}}+\abs{M_{12}},\abs{M_{21}}+\abs{M_{22}}\},
  \]
  and in particular $\abs{\operatorname{diag}(M_1,M_2)}=\max\{\abs{M_1},\abs{M_2}\}$.
\item $\Omega_0^\top=-\Omega_0=\Omega_0^{-1}$, $\Omega_0^2=-I_{2n}$, $\Omega_0^\top\Omega_0=I_{2n}$ and
  $\abs{\Omega_0}=1$. With respect to the splitting $\C^{2n}=\C^n\times\C^n$,
  $\Omega_0\left(\begin{smallmatrix}U_1\\U_2\end{smallmatrix}\right)
  =\left(\begin{smallmatrix}-U_2\\U_1\end{smallmatrix}\right)$ for matrices $U_1,U_2$ with $n$ rows.
\end{enumerate}
\source{HL19 \S2.1, p.~5, and \S3.1, pp.~14--15 (matrix norms and transposes). The transpose bound and
(b) are Lemma~\ref{bp:p1:matrix-norms}.}
\end{lemma}

\begin{proof}
\proofstep{Operator norm}\label{step:p2:matrix-op}
For $\abs v\le1$, $\abs{(Mv)_i}\le\sum_j\abs{M_{ij}}$; equality is attained at
$v_j=\overline{M_{i_0j}}/\abs{M_{i_0j}}$ (and $v_j=0$ where $M_{i_0j}=0$) for a maximizing row $i_0$. Hence
$\abs M=\sup_{\abs v\le1}\abs{Mv}$, and $\abs{M_{ij}}\le\sum_{j'}\abs{M_{ij'}}\le\abs M$.

\proofstep{Transpose and products}\label{step:p2:matrix-transpose}
These are Lemma~\ref{bp:p1:matrix-norms}; directly, $\abs{M^\top}=\max_j\sum_{i=1}^p\abs{M_{ij}}\le
p\max_{i,j}\abs{M_{ij}}\le p\abs M$, and (b) follows from the operator-norm form.

\proofstep{Blocks and $\Omega_0$}\label{step:p2:matrix-blocks}
In (c) the row sums of a block row are the sums of the row sums of its blocks. (d) is a direct
computation from $\Omega_0(a,b)=(-b,a)$; each row of $\Omega_0$ has a single entry $\pm1$.
\end{proof}

\begin{lemma}[Cauchy estimates used below]\label{bp:p2:cauchy}
Let $f$ be holomorphic on $T_r$ with $\abs f\le F$ on $T_r$, and $0<\delta<r$.
\begin{enumerate}[label=(\alph*)]
\item If $f$ is $\C^p$-valued, then $\abs{Df}\le nF/\delta$ and $\abs{(Df)^\top}\le pF/\delta$ on
  $T_{r-\delta}$.
\item If $f$ is vector- or matrix-valued, then $\abs{\partial_\omega f}\le n\abs\omega F/\delta$ on
  $T_{r-\delta}$.
\end{enumerate}
\source{HL19 \S3.1, p.~15.}
\end{lemma}

\begin{proof}
\proofstep{Entries}\label{step:p2:cauchy-entries}
By Lemma~\ref{bp:p1:cauchy-strip}, $\abs{\partial_{\theta_\ell}f}\le F/\delta$ on $T_{r-\delta}$ for each
$\ell$, with the norm of $f$'s values; in particular $\abs{\partial_{\theta_\ell}f_i}\le F/\delta$.

\proofstep{Sums}\label{step:p2:cauchy-sums}
$Df$ is $p\times n$ with entries $\partial_{\theta_\ell}f_i$: its row sums have $n$ terms and its column
sums $p$ terms, which gives (a). For (b), $\abs{\partial_\omega f}\le\sum_\ell\abs{\omega_\ell}
\abs{\partial_{\theta_\ell}f}\le n\abs\omega F/\delta$.
\end{proof}

\begin{lemma}[Averages]\label{bp:p2:averages}
\begin{enumerate}[label=(\alph*)]
\item Let $f:\R^n\to\C^{p\times q}$ be continuous and $1$-periodic. Then $f$ is integrable on $[0,1]^n$,
  $\avg f=\int_{[0,1]^n}f$ (entrywise) is linear in $f$, real if $f$ is real, and
  $\abs{\avg f}\le\sup_{x\in[0,1]^n}\abs{f(x)}$.
\item Let $v$ be holomorphic and $1$-periodic on $T_r$ (not necessarily bounded), with values in
  $\C$, $\C^p$ or $\C^{p\times q}$. Then $\avg{\partial_{\theta_\ell}v}=0$ for every $\ell$ and
  $\avg{\partial_\omega v}=0$.
\end{enumerate}
\source{(a) extends the bound of Lemma~\ref{bp:p1:averages} to continuous matrix-valued functions on
$\R^n$; (b) is Lemma~\ref{bp:p1:averages-real} applied entrywise.}
\end{lemma}

\begin{proof}
\proofstep{Integrability and the bound}\label{step:p2:avg-bound}
A continuous function is bounded and integrable on the compact box $[0,1]^n$, which has measure~$1$.
Then $\abs{\avg f}=\max_i\sum_j\bigl|\int f_{ij}\bigr|\le\max_i\int\sum_j\abs{f_{ij}}\le\sup\abs f$.
Linearity and realness are properties of the integral.

\proofstep{Derivatives}\label{step:p2:avg-deriv}
The restriction of each entry of $v$ to $\R^n\subset T_r$ is continuous, $1$-periodic and
continuously differentiable, with partial derivatives the restrictions of $\partial_{\theta_\ell}v$.
Lemma~\ref{bp:p1:averages-real} gives (b). (Alternatively: $v$ is bounded on $T_{r/2}$ by
Lemma~\ref{bp:p1:periodic-bounded}, and Lemma~\ref{bp:p1:averages} applies there.)
\end{proof}

The cohomological equation of Part~\ref{bp:p1} is stated componentwise. The next lemma restates it in
the operator norms used here, and adds the a-priori form used for the Lagrangian defect.

\begin{lemma}[Cohomological equation in operator norms]\label{bp:p2:cohomological}
Let $0<\delta<r$ and write $\Gamma_\delta:=c_R\gamma^{-1}\delta^{-\varkappa}$.
\begin{enumerate}[label=(\alph*)]
\item\emph{(Solution.)} Let $g:T_r\to\C^p$ be holomorphic and $1$-periodic, with $\abs g\le H$ on $T_r$
  and $\avg g=0$. There is a unique $u=\mathcal R g$, holomorphic and $1$-periodic on $T_{r-\delta}$, with
  $\avg u=0$ and $\partial_\omega u=g$ on $T_{r-\delta}$. It satisfies $\abs u\le\Gamma_\delta H$ on
  $T_{r-\delta}$, and $u$ is real on reals if $g$ is.
\item\emph{(A-priori bound.)} Let $u:T_r\to\C^{p\times q}$ be holomorphic and $1$-periodic with
  $\avg u=0$ and $\abs{\partial_\omega u}\le H$ on $T_r$. Then $\abs u\le\Gamma_\delta H$ on
  $T_{r-\delta}$.
\end{enumerate}
\source{HL19 Lemma~4.1, p.~20, with exponent $\tau+n$ instead of $\tau$ (decision D12).}
\end{lemma}

\begin{proof}
\proofstep{Componentwise solution}\label{step:p2:coh-components}
Apply Lemma~\ref{bp:p1:cohomological} to each component $g_i$, which is holomorphic, periodic, bounded
by $H$ and has zero average. The unique solutions $u_i$ give $u=(u_i)$; uniqueness of $u$ follows from
uniqueness of each $u_i$. Each $\abs{u_i}\le c(n,\tau)\gamma^{-1}\delta^{-\varkappa}H\le\Gamma_\delta H$,
and the max norm gives $\abs u\le\Gamma_\delta H$. Realness is part of Lemma~\ref{bp:p1:cohomological}.

\proofstep{Reduction of (b) to vectors}\label{step:p2:coh-apriori}
Fix $v\in\C^q$ with $\abs v\le1$. The function $uv$ is holomorphic and periodic on $T_r$, has
$\avg{uv}=\avg uv=0$, and $\partial_\omega(uv)=(\partial_\omega u)v=:g$ satisfies $\abs g\le H$ on $T_r$
and $\avg g=0$ (Lemma~\ref{bp:p2:averages}(b)). By the uniqueness in (a), $uv=\mathcal Rg$ on
$T_{r-\delta}$, so $\abs{u(\theta)v}\le\Gamma_\delta H$ there. Taking the supremum over $v$
(Lemma~\ref{bp:p2:matrix}(a)) gives $\abs{u(\theta)}\le\Gamma_\delta H$.
\end{proof}

\begin{lemma}[The Hamiltonian vector field]\label{bp:p2:hamiltonian}
Under Assumption~\ref{bp:p2:standing}:
\begin{enumerate}[label=(\alph*)]
\item $X_h$ is holomorphic on $\mathfrak B$; $DX_h$ and $D^2X_h$ exist there, and
  $\abs{D(DX_h)(z)}=\abs{D^2X_h(z)}\le M$.
\item $X_h(z+(k,0))=X_h(z)$ and $DX_h(z+(k,0))=DX_h(z)$ for $z\in\mathfrak B$, $k\in\Z^n$.
\item $X_h(z)$ and $DX_h(z)$ are real for real $z\in\mathfrak B$.
\item $\Omega_0X_h=(Dh)^\top=:\nabla h$ on $\mathfrak B$, the matrix $\Omega_0DX_h(z)$ is symmetric, and
  \[
  DX_h(z)^\top\Omega_0+\Omega_0DX_h(z)=0\qquad(z\in\mathfrak B).
  \]
\item If $\mathbb B_\beta(z)\subset\mathfrak B$ and $\abs w<\beta$, then
  $\abs{DX_h(z+w)-DX_h(z)}\le M\abs w$ and $\abs{X_h(z+w)-X_h(z)-DX_h(z)w}\le M\abs w^2$.
\end{enumerate}
\source{HL19 Rem.~2.1, p.~6 ($\Omega X_h=(Dh)^\top$); (4.12), p.~23 (canonical case); Rem.~2.3, p.~7
(periodicity).}
\end{lemma}

\begin{proof}
\proofstep{Regularity}\label{step:p2:ham-regular}
A function holomorphic on an open set of $\C^{2n}$ is analytic there, so its derivatives of every
order are holomorphic. By definition $X_h=(\partial_Ih,-\partial_\theta h)$ is built from $Dh$, hence
holomorphic on $\mathfrak B$. The derivative of $DX_h$ is $D^2X_h$ under the canonical isometry
between $\mathcal L(V,\mathcal L(V,W))$ and bilinear maps, so the norms agree.

\proofstep{Periodicity}\label{step:p2:ham-periodic}
The translation $t_k(\theta,I)=(\theta+k,I)$ maps the open set $\mathfrak B$ into itself and
$h\circ t_k=h$ there. Differentiating, $Dh(t_kz)=Dh(z)$, so $X_h\circ t_k=X_h$; differentiating again,
$DX_h\circ t_k=DX_h$.

\proofstep{Realness}\label{step:p2:ham-real}
The real points of $\mathfrak B$ form an open subset of $\R^{2n}$ on which $h$ is real. A complex
partial derivative at a real point is a limit of real difference quotients along a real direction, so
$Dh$, hence $X_h$, is real at real points; the same argument applied to $X_h$ gives $DX_h$.

\proofstep{Symplectic identity}\label{step:p2:ham-symplectic}
With $\Omega_0(a,b)=(-b,a)$, $\Omega_0X_h=\Omega_0(\partial_Ih,-\partial_\theta h)=(\partial_\theta
h,\partial_Ih)=(Dh)^\top$. Differentiating, $\Omega_0DX_h=D^2h$, the Hessian, which is symmetric (Schwarz).
Hence $DX_h^\top\Omega_0=-(\Omega_0DX_h)^\top=-\Omega_0DX_h$, using $\Omega_0^\top=-\Omega_0$.

\proofstep{Mean value estimates}\label{step:p2:ham-mvt}
The ball $\mathbb B_\beta(z)$ is convex and contained in $\mathfrak B$, so it contains the segment
$[z,z+w]$. The mean value inequality for $DX_h$ on this segment, with $\abs{D(DX_h)}\le M$ by (a), gives
the first bound. Lemma~\ref{bp:p1:taylor} applied to $F=X_h$ on $U=\mathfrak B$ with $A=M$ gives
$\abs{X_h(z+w)-X_h(z)-DX_h(z)w}\le\frac M2\abs w^2\le M\abs w^2$.
\end{proof}

\begin{remark}[Transposed derivatives]\label{bp:p2:transposes}
HL19 assume separate bounds $\sigma_{K^\top}$ on $(DK)^\top$ and $c_{X_h^\top,1}$ on $(DX_h)^\top$. Here
$\abs{(DK)^\top}\le2n\abs{DK}$ by Lemma~\ref{bp:p2:matrix}(a), and in the canonical case
$(DX_h)^\top=\Omega_0DX_h\Omega_0$ by Lemma~\ref{bp:p2:hamiltonian}(d), so $\abs{(DX_h)^\top}=
\abs{DX_h}$. The latter is not needed below.
\end{remark}

%----------------------------------------------------------------------------------------------------
\subsection{The frame and the torsion of a parameterization}\label{bp:p2:sec-objects}

\begin{definition}[Objects attached to a parameterization]\label{bp:p2:objects}
Let $r>0$ and let $K:\C^n\to\C^n\times\C^n$ be holomorphic on $T_r$ with $K(T_r)\subset\mathfrak B$.
On $T_r$ define the $2n\times n$ tangent frame, the Gram matrix, the invariance error, the linearized
field along $K$ and the pull-back of the symplectic form:
\[
L:=DK,\qquad G:=L^\top L,\qquad E:=X_h\circ K-L\omega,\qquad A:=DX_h\circ K,\qquad
\Omega_K:=L^\top\Omega_0L.
\]
Where $\det G\ne0$ on $T_r$, define further
\begin{gather*}
B:=G^{-1},\qquad N:=\Omega_0LB,\qquad P:=(L\ \ N),\qquad X_N:=AN-\partial_\omega N,\qquad
T:=N^\top\Omega_0X_N,\\
E_{\mathrm{sym}}:=P^\top\Omega_0P-\Omega_0,\qquad Q:=-\Omega_0P^\top\Omega_0,\qquad
\eta=\begin{pmatrix}\eta^L\\\eta^N\end{pmatrix}:=\Omega_0P^\top\Omega_0E
=\begin{pmatrix}-N^\top\Omega_0E\\ L^\top\Omega_0E\end{pmatrix},
\end{gather*}
and, when $T$ is continuous on $\R^n$, $\avg T:=\int_{[0,1]^n}T(x)\,dx$. Objects of another
parameterization $\tilde K$ (or $\bar K$, $K_s$) carry the same decoration: $\tilde L,\tilde B,\tilde
N,\tilde T,\dots$
\frozen{KAM.gramMatrix, KAM.normalFrame, KAM.torsion,}
\frozen{KAM.torsionAverage, KAM.invarianceError}
\source{HL19 (2.2), (2.12), (3.2)--(3.9), pp.~7, 9, 16--17; (4.29), p.~25; (5.6), p.~32. Canonical
Case~III: Rem.~2.5, p.~12, and Rem.~3.3, p.~18.}
\end{definition}

\begin{remark}[Agreement with the frozen definitions]\label{bp:p2:frozen-agreement}
On the open strip $T_r$, the Fr\'echet derivative of $K$ is $DK$, and the frozen \Lean{gramMatrix},
\Lean{normalFrame} (column $j$ is $\sum_kB_{kj}\Omega_0\partial_{\theta_k}K$, i.e.\ column $j$ of
$\Omega_0LB$), \Lean{torsion} (entry $(i,j)$ is $N_i^\top\Omega_0(AN_j-DN_j\,\omega)$) and
\Lean{invarianceError} coincide with $G$, $N$, $T$ and $E$. The frozen matrix inverse is the true inverse
where the determinant is a unit, and the frozen \Lean{torsionAverage} is the entrywise integral of $T$
over $[0,1]^n$ at real points, i.e.\ $\avg T$. Following Remark~\ref{bp:p0:totality}, every use of $B$,
$N$ or $T$ below is preceded by a proof that $\det G\ne0$, and every use of $\avg T$ by a proof that $T$
is continuous on $\R^n$.
\end{remark}

\begin{definition}[Admissible parameterization]\label{bp:p2:admissible}
Let $r>0$, $\beta>0$ and $0\le\varepsilon\le1$. A map $K:\C^n\to\C^n\times\C^n$ is
\emph{$(r,\beta,\varepsilon)$-admissible} if
\begin{enumerate}[label=(A\arabic*)]
\item\label{bp:p2:A1} $K$ is holomorphic on $T_r$, primary on $T_r$, and real on reals;
\item\label{bp:p2:A2} $\mathbb B_\beta(K(\theta))\subset\mathfrak B$ for every $\theta\in T_r$;
\item\label{bp:p2:A3} $\abs{DK}\le\sigma_K$ on $T_r$;
\item\label{bp:p2:A4} $\det G\ne0$ and $\abs B\le\sigma_B$ on $T_r$;
\item\label{bp:p2:A5} $\avg T$ is invertible and $\abs{\avg T^{-1}}\le\sigma_T$;
\item\label{bp:p2:A6} $\abs E\le\varepsilon$ on $T_r$.
\end{enumerate}
\end{definition}

Condition~\ref{bp:p2:A5} is meaningful because of \ref{bp:p2:A4} and
Lemma~\ref{bp:p2:frame-basic}(a). The hypotheses of Theorem~\ref{bp:p0:g1} state that $K$ is
$(\rho,\mu,e)$-admissible, with the sharper bounds $\sigma_K-\mu$, $\sigma_B-\mu$, $\sigma_T-\mu$ in
\ref{bp:p2:A3}--\ref{bp:p2:A5}, once $e\le1$ is known.

\begin{lemma}[Periodicity and realness of the attached objects]\label{bp:p2:periodicity}
Let $K$ be holomorphic and primary on $T_r$, with $K(T_r)\subset\mathfrak B$. Then $L$, $G$, $E$, $A$,
$\Omega_K$ and $h\circ K$ are $1$-periodic on $T_r$, as are the action components $K^I$ and the functions
$f_j:=(K^I)^\top\partial_{\theta_j}K^\theta$, $j=1,\dots,n$, where $K=(K^\theta,K^I)$. If $\det G\ne0$ on
$T_r$, then $B$, $N$, $P$, $X_N$, $T$, $E_{\mathrm{sym}}$ and $\eta$ are $1$-periodic on $T_r$. If in
addition $K$ is real on reals, all these objects are real at real points, and so is $\avg T$.
\source{HL19 Rem.~2.3, p.~7 (``all expressions and formulas presented in this paper remain valid'').}
\sourcegap{HL19 Rem.~2.3 asserts the periodicity of all objects for primary tori without checking it;
the check is supplied here.}
\end{lemma}

\begin{proof}
\proofstep{Periodicity of $DK$ and $K^I$}\label{step:p2:per-dk}
For $k\in\Z^n$ the identity $K(\theta+k)=K(\theta)+(k,0)$ holds on the open set $T_r$, which is
invariant under $\theta\mapsto\theta+k$. Its second component says $K^I(\theta+k)=K^I(\theta)$;
differentiating it gives $DK(\theta+k)=DK(\theta)$, and differentiating again gives the same for second
derivatives.

\proofstep{Periodicity of the composed objects}\label{step:p2:per-compose}
By Lemma~\ref{bp:p2:hamiltonian}(b) and the periodicity of $h$, $X_h(K(\theta+k))=X_h(K(\theta)+(k,0))
=X_h(K(\theta))$, and likewise for $DX_h\circ K$ and $h\circ K$. Every other object in the list is an
algebraic expression (products, transposes, the matrix inverse) in $L$, $X_h\circ K$, $A$, $K^I$,
derivatives of these, and the constants $\Omega_0,\omega$; derivatives of periodic functions are periodic.

\proofstep{Realness}\label{step:p2:per-real}
If $K$ is real on $\R^n$, then $DK$ and its derivatives are real at real points (real difference
quotients), $K(x)$ is a real point of $\mathfrak B$ for real $x$, and $X_h(K(x))$, $DX_h(K(x))$ are real
by Lemma~\ref{bp:p2:hamiltonian}(c). The inverse of a real invertible matrix is real (adjugate formula),
and $\partial_\omega$ of a function real on $\R^n$ is real on $\R^n$ because $\omega\in\R^n$. Hence all
objects are real at real points, and $\avg T$ is real by Lemma~\ref{bp:p2:averages}(a).
\end{proof}

\begin{lemma}[Frequency bound]\label{bp:p2:frequency}
Let $K$ satisfy \ref{bp:p2:A3}, \ref{bp:p2:A4}, \ref{bp:p2:A6} on $T_r$ with $\varepsilon\le1$, and
$K(T_r)\subset\mathfrak B$. Then
\[
\gamma\le\abs\omega\le2n\sigma_K\sigma_B(M+\varepsilon)\le\Theta.
\]
\source{Ours. The analogous remark for G2 is in ledger D-G1-1 (``invariance with H2--H3 bounds
$\abs\omega$, hence $\gamma$'').}
\end{lemma}

\begin{proof}
\proofstep{Solving for $\omega$}\label{step:p2:freq-solve}
Fix $\theta\in T_r$. From $E=X_h(K)-L\omega$ and $BL^\top L=I_n$ we get
$\omega=B(\theta)L(\theta)^\top\bigl(X_h(K(\theta))-E(\theta)\bigr)$.

\proofstep{Upper bound}\label{step:p2:freq-upper}
By Lemma~\ref{bp:p2:matrix}, $\abs{L^\top}\le2n\abs L\le2n\sigma_K$; with $\abs B\le\sigma_B$,
$\abs{X_h}\le M$ and $\abs E\le\varepsilon$ this gives $\abs\omega\le2n\sigma_K\sigma_B(M+\varepsilon)
\le2n\mathfrak m^2\cdot2\mathfrak m=\Theta$, since $M+\varepsilon\le\mathfrak m+1\le2\mathfrak m$.

\proofstep{Lower bound}\label{step:p2:freq-lower}
The Diophantine condition at $k=e_1$, where $\abs{e_1}_1^\tau=1$, reads $\gamma\le\abs{\omega_1}\le
\abs\omega$.
\end{proof}

\begin{lemma}[Regularity and bounds of the frame]\label{bp:p2:frame-basic}
Let $K$ be holomorphic on $T_r$ and satisfy \ref{bp:p2:A2}--\ref{bp:p2:A4} (primary and real are not
needed here).
\begin{enumerate}[label=(\alph*)]
\item $L$, $G$, $E$, $A$, $\Omega_K$, $B$, $N$, $P$, $X_N$, $T$, $E_{\mathrm{sym}}$, $Q$, $\eta$ are
  holomorphic on $T_r$; $B$ is symmetric; $T$ is continuous on $\R^n$, so $\avg T$ is defined.
\item On $T_r$: $\abs L\le\mathfrak m$, $\abs{L^\top}\le2n\mathfrak m$, $\abs B\le\mathfrak m$,
  $\abs N\le\mathfrak m^2$, $\abs{N^\top}\le2n\mathfrak m^2$, $\abs P\le2\mathfrak m^2$,
  $\abs A\le\mathfrak m$.
\item If \ref{bp:p2:A6} holds with $\varepsilon\le1$, then on $T_r$: $\abs{\eta^L}\le2n\mathfrak
  m^2\varepsilon$, $\abs{\eta^N}\le2n\mathfrak m\varepsilon$.
\item If $\abs\omega\le\Theta$ and $0<\delta\le1$, $\delta<r$, then $\abs{X_N}\le C_X/\delta$ and
  $\abs T\le C_T/\delta$ on $T_{r-\delta}$.
\end{enumerate}
\source{HL19 Lemma~4.5 (4.23)--(4.24), p.~24; Lemma~4.6 (4.31)--(4.35), p.~25; Lemma~4.8
(4.59)--(4.60), pp.~27--29 (by a cruder route, see Remark~\ref{bp:p2:torsion-route}).}
\end{lemma}

\begin{proof}
\proofstep{Holomorphy}\label{step:p2:fb-holo}
$L=DK$ is holomorphic on $T_r$; $X_h\circ K$ and $A=DX_h\circ K$ are holomorphic because $K(T_r)\subset
\mathfrak B$ and Lemma~\ref{bp:p2:hamiltonian}(a). By \ref{bp:p2:A4}, $\det G\ne0$ on $T_r$, so
$B=G^{-1}$ is holomorphic there (Lemma~\ref{bp:p1:holomorphic-inverse}). Then $N,P,\partial_\omega
N,X_N,T,E_{\mathrm{sym}},Q,\eta$ are holomorphic as products and derivatives of holomorphic functions.
$G$ is symmetric, hence so is $B$. A holomorphic function is continuous, and $\R^n\subset T_r$.

\proofstep{Frame bounds}\label{step:p2:fb-bounds}
By \ref{bp:p2:A3}, \ref{bp:p2:A4} and Lemma~\ref{bp:p2:matrix}: $\abs L\le\sigma_K\le\mathfrak m$,
$\abs{L^\top}\le2n\mathfrak m$, $\abs N\le\abs{\Omega_0}\abs L\abs B\le\mathfrak m^2$, $\abs{N^\top}\le
2n\abs N$, $\abs P\le\abs L+\abs N\le2\mathfrak m^2$, and $\abs A\le M\le\mathfrak m$ by
Assumption~\ref{bp:p2:standing}.

\proofstep{Bounds on $\eta$}\label{step:p2:fb-eta}
$\abs{\eta^L}\le\abs{N^\top}\abs{\Omega_0}\abs E\le2n\mathfrak m^2\varepsilon$ and
$\abs{\eta^N}\le\abs{L^\top}\abs E\le2n\mathfrak m\varepsilon$.

\proofstep{Torsion bound}\label{step:p2:fb-torsion}
By Lemma~\ref{bp:p2:cauchy}(b) applied to $N$ ($\abs N\le\mathfrak m^2$ on $T_r$),
$\abs{\partial_\omega N}\le n\Theta\mathfrak m^2/\delta=4n^2\mathfrak m^5/\delta$ on $T_{r-\delta}$.
Hence $\abs{X_N}\le\abs A\abs N+\abs{\partial_\omega N}\le\mathfrak m^3+4n^2\mathfrak m^5/\delta\le
5n^2\mathfrak m^5/\delta=C_X/\delta$, using $\delta\le1$. Then $\abs T\le\abs{N^\top}\abs{X_N}\le
2n\mathfrak m^2C_X/\delta=C_T/\delta$.
\end{proof}

\begin{remark}[Route for the torsion]\label{bp:p2:torsion-route}
HL19 bound $T$ without a factor $1/\delta$ by expanding $L_\omega N$ through $L_\omega K=E-X_h\circ K$
(Lemma~4.8 and condition (4.40), $\norm E_\rho/\delta<\nu$). We use instead the Cauchy estimate and the
frequency bound of Lemma~\ref{bp:p2:frequency}. This costs a factor $1/\delta$ in the torsion, and later
one more strip loss, but no power of $\gamma$.
\end{remark}

%----------------------------------------------------------------------------------------------------
\subsection{Geometry of an approximately invariant torus}\label{bp:p2:sec-geometry}

The identities of this subsection are exact on $T_r$; they express how far the frame $P$ is from being
symplectic and from reducing the linearized field, in terms of $\Omega_K$ and $DE$. The next
subsection then estimates $\Omega_K$ through the cohomological equation.

\begin{lemma}[Frame identities]\label{bp:p2:frame-identities}
Let $K$ be holomorphic on $T_r$, with $K(T_r)\subset\mathfrak B$ and $\det G\ne0$ on $T_r$. Then on
$T_r$:
\begin{enumerate}[label=(\alph*)]
\item $L^\top\Omega_0N=-I_n$, $N^\top\Omega_0L=I_n$ and $N^\top\Omega_0N=B\Omega_KB$;
\item $P^\top\Omega_0P=\Omega_0+E_{\mathrm{sym}}$ with $E_{\mathrm{sym}}=\operatorname{diag}(\Omega_K,
  B\Omega_KB)$;
\item $\partial_\omega L=AL-DE$;
\item $\partial_\omega\Omega_K=-(DE)^\top\Omega_0L-L^\top\Omega_0DE$;
\item $Q\,(AP-\partial_\omega P)=\Lambda+E_{\mathrm{red}}$, where
  \[
  \Lambda=\begin{pmatrix}0&T\\0&0\end{pmatrix},\qquad
  E_{\mathrm{red}}=\begin{pmatrix}N^\top\Omega_0DE&0\\-L^\top\Omega_0DE&(DE)^\top\Omega_0N\end{pmatrix};
  \]
\item $QP=I_{2n}-\Omega_0E_{\mathrm{sym}}$ and $QE=-\eta$.
\end{enumerate}
\source{HL19 (4.36)--(4.38), p.~26, and Rem.~4.7, p.~26 (a, b); (4.13) and the line after it, p.~23
(c, d); Lemma~4.10, (4.61)--(4.66), pp.~29--30 (e); (5.4), p.~31 (f). Cross-check: GJLV05 Lemma~20,
p.~891.}
\sourcegap{HL19 Rem.~4.7 leaves the Case~III identities (a)--(b) ``as an exercise''; they are proved
here. In the canonical Case~III, (e) is sharper than HL19 (4.66): the $(2,2)$ block reduces exactly to
$(DE)^\top\Omega_0N$.}
\end{lemma}

\begin{proof}
\proofstep{Products with $\Omega_0$}\label{step:p2:fi-products}
Using $\Omega_0^2=-I_{2n}$, $\Omega_0^\top=-\Omega_0$, $B^\top=B$ and $L^\top L=G=B^{-1}$:
$L^\top\Omega_0N=L^\top\Omega_0^2LB=-GB=-I_n$; $N^\top\Omega_0L=-BL^\top\Omega_0^2L=BG=I_n$;
$N^\top\Omega_0N=-BL^\top\Omega_0^3LB=BL^\top\Omega_0LB=B\Omega_KB$. This is (a).

\proofstep{Symplecticity defect}\label{step:p2:fi-sym}
In blocks, $P^\top\Omega_0P=\left(\begin{smallmatrix}L^\top\Omega_0L&L^\top\Omega_0N\\
N^\top\Omega_0L&N^\top\Omega_0N\end{smallmatrix}\right)=\left(\begin{smallmatrix}\Omega_K&-I_n\\
I_n&B\Omega_KB\end{smallmatrix}\right)$ by (a), and subtracting
$\Omega_0=\left(\begin{smallmatrix}0&-I_n\\I_n&0\end{smallmatrix}\right)$ gives (b).

\proofstep{Derivative of the frame along $\omega$}\label{step:p2:fi-dl}
$K$ is holomorphic, so mixed partial derivatives commute and
$\partial_\omega\partial_{\theta_j}K=\partial_{\theta_j}(L\omega)$. Since $L\omega=X_h\circ K-E$ on $T_r$,
the chain rule gives $\partial_{\theta_j}(L\omega)=A\,\partial_{\theta_j}K-\partial_{\theta_j}E$, which
is column $j$ of (c).

\proofstep{Derivative of $\Omega_K$}\label{step:p2:fi-domegak}
By the product rule and (c),
$\partial_\omega\Omega_K=(AL-DE)^\top\Omega_0L+L^\top\Omega_0(AL-DE)
=L^\top(A^\top\Omega_0+\Omega_0A)L-(DE)^\top\Omega_0L-L^\top\Omega_0DE$, and the first term vanishes
by Lemma~\ref{bp:p2:hamiltonian}(d) at the points $K(\theta)\in\mathfrak B$.

\proofstep{Block form of the reduced field}\label{step:p2:fi-blocks}
Put $X_L:=AL-\partial_\omega L$ and recall $X_N=AN-\partial_\omega N$, so $AP-\partial_\omega
P=(X_L\ X_N)$. By Lemma~\ref{bp:p2:matrix}(d),
\[
Q(X_L\ X_N)=-\Omega_0\begin{pmatrix}L^\top\Omega_0X_L&L^\top\Omega_0X_N\\N^\top\Omega_0X_L&N^\top\Omega_0X_N
\end{pmatrix}=\begin{pmatrix}N^\top\Omega_0X_L&N^\top\Omega_0X_N\\-L^\top\Omega_0X_L&-L^\top\Omega_0X_N
\end{pmatrix}.
\]
By (c), $X_L=DE$, and $N^\top\Omega_0X_N=T$ by definition.

\proofstep{The $(2,2)$ block}\label{step:p2:fi-22}
Differentiating the constant identity $L^\top\Omega_0N=-I_n$ of (a) along $\omega$ gives
$(\partial_\omega L)^\top\Omega_0N+L^\top\Omega_0\partial_\omega N=0$. Hence, with (c) and
Lemma~\ref{bp:p2:hamiltonian}(d),
\[
L^\top\Omega_0X_N=L^\top\Omega_0AN+(AL-DE)^\top\Omega_0N=L^\top(\Omega_0A+A^\top\Omega_0)N-(DE)^\top
\Omega_0N=-(DE)^\top\Omega_0N.
\]
Together with the previous step this is (e).

\proofstep{The products with $Q$}\label{step:p2:fi-q}
$QP=-\Omega_0(\Omega_0+E_{\mathrm{sym}})=I_{2n}-\Omega_0E_{\mathrm{sym}}$ by (b), and $QE=-\eta$ by the
definition of $\eta$. The block form of $\eta$ in Definition~\ref{bp:p2:objects} is
Lemma~\ref{bp:p2:matrix}(d) applied to $P^\top\Omega_0E=(L^\top\Omega_0E;N^\top\Omega_0E)$.
\end{proof}

\begin{lemma}[Exactness: two vanishing averages]\label{bp:p2:exactness}
Let $K$ be holomorphic and primary on $T_r$ with $K(T_r)\subset\mathfrak B$. Then
\[
\avg{\Omega_K}=0\qquad\text{and}\qquad\avg{L^\top\Omega_0E}=0,\ \text{ i.e. }\ \avg{\eta^N}=0
\ \text{ when }\det G\ne0.
\]
\source{HL19 Lemma~4.3, (4.10), p.~22; Lemma~4.5, (4.16) and its proof, p.~24. Cross-check: GJLV05
Lemma~19(a), p.~889.}
\end{lemma}

\begin{proof}
\proofstep{$\Omega_K$ is a curl of periodic functions}\label{step:p2:ex-curl}
Write $K=(K^\theta,K^I)$. Since $u^\top\Omega_0v=u^I\cdot v^\theta-u^\theta\cdot v^I$ for
$u=(u^\theta,u^I)$, $v=(v^\theta,v^I)$, the entries of $\Omega_K$ are
$(\Omega_K)_{ij}=\partial_{\theta_i}K^I\cdot\partial_{\theta_j}K^\theta-\partial_{\theta_i}K^\theta\cdot
\partial_{\theta_j}K^I$. With $f_j=(K^I)^\top\partial_{\theta_j}K^\theta$ and the symmetry of second
derivatives, $(\Omega_K)_{ij}=\partial_{\theta_i}f_j-\partial_{\theta_j}f_i$. This is the pull-back of
the primitive $\alpha_0=\sum_iI_i\,d\theta_i$, whose coefficient $K^I$ is periodic for a primary torus.

\proofstep{Average of $\Omega_K$}\label{step:p2:ex-avg-omega}
The $f_j$ are holomorphic and periodic on $T_r$ (Lemma~\ref{bp:p2:periodicity}), so
$\avg{\partial_{\theta_i}f_j}=0$ by Lemma~\ref{bp:p2:averages}(b), and $\avg{\Omega_K}=0$.

\proofstep{The gradient identity}\label{step:p2:ex-gradient}
By Lemma~\ref{bp:p2:hamiltonian}(d), $\Omega_0X_h=(Dh)^\top$, so
$L^\top\Omega_0X_h(K)=(Dh(K)L)^\top=(D(h\circ K))^\top$, and
$L^\top\Omega_0E=(D(h\circ K))^\top-\Omega_K\omega$.

\proofstep{Average of $L^\top\Omega_0E$}\label{step:p2:ex-avg-eta}
$h\circ K$ is holomorphic and periodic on $T_r$ (Lemma~\ref{bp:p2:periodicity}), so its partial
derivatives have zero average (Lemma~\ref{bp:p2:averages}(b)); and $\avg{\Omega_K\omega}=\avg{\Omega_K}
\omega=0$. Hence $\avg{L^\top\Omega_0E}=0$, and $\eta^N=L^\top\Omega_0E$.
\end{proof}

\begin{lemma}[Approximate Lagrangian property, symplecticity and reducibility]\label{bp:p2:geometry}
Let $K$ be $(r,\beta,\varepsilon)$-admissible ($\varepsilon\le1$), and let $0<\delta\le1$, $2\delta<r$.
Then
\begin{align*}
\abs{DE}&\le n\varepsilon/\delta,\quad\abs{(DE)^\top}\le2n\varepsilon/\delta,\quad
\abs{\partial_\omega\Omega_K}\le2n(n+1)\mathfrak m\,\varepsilon/\delta,\quad
\abs{E_{\mathrm{red}}}\le C_{\mathrm{red}}\,\varepsilon/\delta&&\text{on }T_{r-\delta},\\
\abs{\Omega_K}&\le\frac{C_\Omega\,\varepsilon}{\gamma\delta^{\varkappa+1}},\qquad
\abs{E_{\mathrm{sym}}}\le\frac{C_{\mathrm{sym}}\,\varepsilon}{\gamma\delta^{\varkappa+1}}&&\text{on }
T_{r-2\delta}.
\end{align*}
\source{HL19 Lemma~4.3, (4.11), (4.14)--(4.15), pp.~22--23; Lemma~4.6, (4.30), p.~25, with Rem.~4.7;
Lemma~4.10, (4.70)--(4.73), p.~30. Cross-check: GJLV05 Lemma~19(d), p.~890, and Lemma~20 (109),
p.~891.}
\end{lemma}

\begin{proof}
\proofstep{Derivatives of the error}\label{step:p2:geo-de}
$E$ is holomorphic on $T_r$ with $\abs E\le\varepsilon$ and has $2n$ components, so
Lemma~\ref{bp:p2:cauchy}(a) gives the bounds on $DE$ and $(DE)^\top$ on $T_{r-\delta}$.

\proofstep{Bound on $\partial_\omega\Omega_K$}\label{step:p2:geo-domegak}
By Lemma~\ref{bp:p2:frame-identities}(d) and Lemma~\ref{bp:p2:frame-basic}(b), on $T_{r-\delta}$:
$\abs{\partial_\omega\Omega_K}\le\abs{(DE)^\top}\abs L+\abs{L^\top}\abs{DE}\le
(2n\mathfrak m+2n^2\mathfrak m)\varepsilon/\delta$.

\proofstep{Bound on $E_{\mathrm{red}}$}\label{step:p2:geo-ered}
By Lemma~\ref{bp:p2:frame-identities}(e) and Lemma~\ref{bp:p2:matrix}(c), on $T_{r-\delta}$ the first
block row has norm $\le\abs{N^\top}\abs{DE}\le2n^2\mathfrak m^2\varepsilon/\delta$ and the second
$\le\abs{L^\top}\abs{DE}+\abs{(DE)^\top}\abs N\le(2n^2\mathfrak m+2n\mathfrak m^2)\varepsilon/\delta$;
both are $\le C_{\mathrm{red}}\varepsilon/\delta$.

\proofstep{The Lagrangian defect by the cohomological equation}\label{step:p2:geo-omegak}
$\Omega_K$ is holomorphic and periodic on $T_{r-\delta}$, with $\avg{\Omega_K}=0$
(Lemma~\ref{bp:p2:exactness}) and $\abs{\partial_\omega\Omega_K}\le2n(n+1)\mathfrak m\varepsilon/\delta$
there. Lemma~\ref{bp:p2:cohomological}(b) on the strip $T_{r-\delta}$ with loss $\delta$ (note
$\delta<r-\delta$) gives $\abs{\Omega_K}\le c_R\gamma^{-1}\delta^{-\varkappa}\cdot2n(n+1)\mathfrak m
\varepsilon/\delta=C_\Omega\varepsilon/(\gamma\delta^{\varkappa+1})$ on $T_{r-2\delta}$.

\proofstep{The symplecticity defect}\label{step:p2:geo-esym}
By Lemma~\ref{bp:p2:frame-identities}(b) and Lemma~\ref{bp:p2:matrix}(c),
\[
\abs{E_{\mathrm{sym}}}=\max\{\abs{\Omega_K},\abs{B\Omega_KB}\}\le\max\{1,\mathfrak m^2\}\abs{\Omega_K}
=\mathfrak m^2\abs{\Omega_K},
\]
which gives the last bound.
\end{proof}

The reducibility defect $E_{\mathrm{red}}$ carries no small divisor: in the canonical Case~III the
Lagrangian defect enters the linearized equation only through $E_{\mathrm{sym}}$, i.e.\ through
$\Omega_K$.

%----------------------------------------------------------------------------------------------------
\subsection{The Newton correction}\label{bp:p2:sec-correction}

The Newton correction $\Delta K$ should solve the linearized equation $A\Delta K-\partial_\omega\Delta
K=-E$ (HL19 (5.1), p.~31, where $L_\omega=-\partial_\omega$). In the frame, $\Delta K=P\xi$, and
multiplying by $Q$ turns the equation, up to the defects $E_{\mathrm{red}}$ and $E_{\mathrm{sym}}$, into the
triangular system $\Lambda\xi-\partial_\omega\xi=\eta$:
\[
\partial_\omega\xi^N=-\eta^N,\qquad\partial_\omega\xi^L=T\xi^N-\eta^L.
\]
The first equation is solvable because $\avg{\eta^N}=0$ (Lemma~\ref{bp:p2:exactness}); its solution is
fixed up to a constant $\xi^N_0$, and this constant is chosen, by inverting $\avg T$, to make the right
side of the second equation have zero average.

\begin{lemma}[The triangular system]\label{bp:p2:triangular}
Let $0<\delta$, $2\delta<r$, and write $\Gamma_\delta=c_R\gamma^{-1}\delta^{-\varkappa}$. Let
$\eta^L,\eta^N:T_r\to\C^n$ be holomorphic and $1$-periodic with $\abs{\eta^L}\le H_L$,
$\abs{\eta^N}\le H_N$ on $T_r$ and $\avg{\eta^N}=0$. Let $T:T_r\to\C^{n\times n}$ be holomorphic and
$1$-periodic, with $\abs T\le H_T$ on $T_{r-\delta}$, $\avg T$ invertible and $\abs{\avg T^{-1}}\le
\sigma$. Define
\begin{gather*}
u^N:=\mathcal R\eta^N\ \text{on }T_{r-\delta},\qquad
\xi^N_0:=\avg T^{-1}\bigl(\avg{\eta^L}+\avg{Tu^N}\bigr)\in\C^n,\qquad\xi^N:=\xi^N_0-u^N,\\
g^L:=T\xi^N-\eta^L\ \text{on }T_{r-\delta},\qquad\xi^L:=\mathcal Rg^L\ \text{on }T_{r-2\delta},\qquad
\xi:=(\xi^L,\xi^N),
\end{gather*}
with $\mathcal R$ from Lemma~\ref{bp:p2:cohomological}(a). Then:
\begin{enumerate}[label=(\alph*)]
\item $\avg{g^L}=0$, so $\xi^L$ is defined; $\xi^N$ and $g^L$ are holomorphic and periodic on
  $T_{r-\delta}$, and $\xi$ on $T_{r-2\delta}$;
\item $\partial_\omega\xi^N=-\eta^N$ on $T_{r-\delta}$ and $\partial_\omega\xi^L=g^L$ on
  $T_{r-2\delta}$; equivalently $\Lambda\xi-\partial_\omega\xi=\eta$ on $T_{r-2\delta}$, with
  $\Lambda=\left(\begin{smallmatrix}0&T\\0&0\end{smallmatrix}\right)$;
\item $\abs{\xi^N_0}\le\sigma(H_L+H_T\Gamma_\delta H_N)$; $\abs{\xi^N}\le\abs{\xi^N_0}+\Gamma_\delta
  H_N$ and $\abs{g^L}\le H_T\sup_{T_{r-\delta}}\abs{\xi^N}+H_L$ on $T_{r-\delta}$;
  $\abs{\xi^L}\le\Gamma_\delta\sup_{T_{r-\delta}}\abs{g^L}$ on $T_{r-2\delta}$;
\item if $\eta^L,\eta^N,T$ are real on reals, then so are $\xi^N_0$, $\xi$ and $g^L$.
\end{enumerate}
\source{HL19 Lemma~5.1, pp.~31--32 (with $\xi^L_0=0$ as chosen on p.~33); signs converted by
$L_\omega=-\partial_\omega$.}
\sourcegap{HL19 do not prove Lemma~5.1 (``directly adapted from Lemma~4.14 in [33]''); the proof is
supplied here.}
\end{lemma}

\begin{proof}
\proofstep{First cohomological equation}\label{step:p2:tri-first}
By Lemma~\ref{bp:p2:cohomological}(a) on $T_r$, $u^N$ is holomorphic, periodic and real (if $\eta^N$
is) on $T_{r-\delta}$, with $\avg{u^N}=0$, $\partial_\omega u^N=\eta^N$ and $\abs{u^N}\le\Gamma_\delta
H_N$. Hence $\partial_\omega\xi^N=-\eta^N$ and $\abs{\xi^N}\le\abs{\xi^N_0}+\Gamma_\delta H_N$.

\proofstep{The torsion average fixes $\xi^N_0$}\label{step:p2:tri-xi0}
$Tu^N$ is continuous on $\R^n$, so $\avg{Tu^N}$ is defined, and $\abs{\avg{Tu^N}}\le H_T\Gamma_\delta
H_N$, $\abs{\avg{\eta^L}}\le H_L$ (Lemma~\ref{bp:p2:averages}(a)). This gives the bound on
$\xi^N_0$. By linearity of the average, $\avg{g^L}=\avg T\xi^N_0-\avg{Tu^N}-\avg{\eta^L}=0$.

\proofstep{Second cohomological equation}\label{step:p2:tri-second}
$g^L$ is holomorphic and periodic on $T_{r-\delta}$, bounded by $H_T\sup\abs{\xi^N}+H_L$, with zero
average. Lemma~\ref{bp:p2:cohomological}(a) on $T_{r-\delta}$ with loss $\delta$ (note $\delta<r-\delta$)
gives $\xi^L$ on $T_{r-2\delta}$ with $\partial_\omega\xi^L=g^L$ and the stated bound.

\proofstep{Triangular form}\label{step:p2:tri-form}
$\Lambda\xi-\partial_\omega\xi=(T\xi^N-\partial_\omega\xi^L,\,-\partial_\omega\xi^N)=(\eta^L,\eta^N)$.

\proofstep{Realness}\label{step:p2:tri-real}
If the data are real on reals, $\avg T$, $\avg{\eta^L}$, $\avg{Tu^N}$ are real, hence $\xi^N_0$ is real,
and Lemma~\ref{bp:p2:cohomological}(a) preserves realness at each stage.
\end{proof}

\begin{lemma}[Size of the correction]\label{bp:p2:correction}
Let $K$ be $(r,\beta,\varepsilon)$-admissible and $0<\delta\le1$, $2\delta<r$. Apply
Lemma~\ref{bp:p2:triangular} to $\eta^L,\eta^N$ and the torsion $T$ of $K$, and put $\Delta K:=P\xi$ on
$T_{r-2\delta}$. Then $\xi$ and $\Delta K$ are holomorphic, $1$-periodic and real on reals, and
\begin{align*}
\abs{\xi^N}&\le\frac{C_{\xi N}\varepsilon}{\gamma\delta^{\varkappa+1}},\quad
\abs{g^L}\le\frac{C_g\varepsilon}{\gamma\delta^{\varkappa+2}}&&\text{on }T_{r-\delta},\\
\abs\xi&\le\frac{C_\xi\varepsilon}{\gamma^2\delta^a},\quad
\abs{\Delta K}\le\frac{C_{\Delta K}\varepsilon}{\gamma^2\delta^a},\quad
\abs{\partial_\omega\xi}\le\frac{C_g\varepsilon}{\gamma\delta^{\varkappa+2}}&&\text{on }T_{r-2\delta}.
\end{align*}
\source{HL19 (5.6), p.~32; (5.20)--(5.24), pp.~33--34; (5.26)--(5.28), pp.~34--35.}
\end{lemma}

\begin{proof}
\proofstep{The data of the triangular system}\label{step:p2:cor-data}
By Lemma~\ref{bp:p2:frequency}, $\abs\omega\le\Theta$ and $\gamma\le\Theta$. By
Lemmas~\ref{bp:p2:periodicity}, \ref{bp:p2:frame-basic} and~\ref{bp:p2:exactness}: $\eta^L,\eta^N,T$ are
holomorphic, periodic and real on reals; $H_L=2n\mathfrak m^2\varepsilon$, $H_N=2n\mathfrak m\varepsilon$,
$H_T=C_T/\delta$; $\avg{\eta^N}=0$; and $\sigma=\sigma_T\le\mathfrak m$ by \ref{bp:p2:A5}.

\proofstep{Absorbing powers of $\gamma\delta$}\label{step:p2:cor-absorb}
Since $\gamma\le\Theta$ and $\delta\le1$, $\gamma\delta^j\le\Theta$ for every $j\ge0$; hence
$\varepsilon\le\Theta\varepsilon/(\gamma\delta^j)$ and $\varepsilon/(\gamma\delta^j)\le\Theta
\varepsilon/(\gamma^2\delta^{j+i})$ for $i,j\ge0$. These inequalities are used in the next two steps.

\proofstep{Bounds on $\xi^N$ and $g^L$}\label{step:p2:cor-xin}
Lemma~\ref{bp:p2:triangular}(c): $\abs{\xi^N_0}\le\mathfrak m(2n\mathfrak m^2\varepsilon+
(C_T/\delta)c_R\gamma^{-1}\delta^{-\varkappa}2n\mathfrak m\varepsilon)\le2n\mathfrak m^2(\mathfrak
m\Theta+c_RC_T)\varepsilon/(\gamma\delta^{\varkappa+1})=C_{\xi_0}\varepsilon/(\gamma\delta^{\varkappa+1})$;
then $\abs{\xi^N}\le(C_{\xi_0}+2nc_R\mathfrak m)\varepsilon/(\gamma\delta^{\varkappa+1})$ and
$\abs{g^L}\le(C_T/\delta)C_{\xi N}\varepsilon/(\gamma\delta^{\varkappa+1})+2n\mathfrak m^2\varepsilon\le
C_g\varepsilon/(\gamma\delta^{\varkappa+2})$ on $T_{r-\delta}$.

\proofstep{Bounds on $\xi$ and $\Delta K$}\label{step:p2:cor-xi}
$\abs{\xi^L}\le c_R\gamma^{-1}\delta^{-\varkappa}C_g\varepsilon/(\gamma\delta^{\varkappa+2})=c_RC_g
\varepsilon/(\gamma^2\delta^a)$ on $T_{r-2\delta}$, and $\abs{\xi^N}\le\Theta C_{\xi N}\varepsilon/
(\gamma^2\delta^a)$ by the absorption step; so $\abs\xi=\max\{\abs{\xi^L},\abs{\xi^N}\}\le C_\xi
\varepsilon/(\gamma^2\delta^a)$, and $\abs{\Delta K}\le\abs P\abs\xi\le2\mathfrak m^2C_\xi\varepsilon/
(\gamma^2\delta^a)$.

\proofstep{Bound on $\partial_\omega\xi$}\label{step:p2:cor-dxi}
By Lemma~\ref{bp:p2:triangular}(b), $\partial_\omega\xi=(g^L,-\eta^N)$, and by the absorption step
\[
\abs{\eta^N}\le2n\mathfrak m\varepsilon\le\frac{2n\mathfrak m^2\Theta\varepsilon}{\gamma
\delta^{\varkappa+2}}\le\frac{C_g\varepsilon}{\gamma\delta^{\varkappa+2}}.
\]

\proofstep{Regularity of the correction}\label{step:p2:cor-reg}
$P$ and $\xi$ are holomorphic, periodic and real on reals (Lemmas~\ref{bp:p2:periodicity},
\ref{bp:p2:frame-basic} and~\ref{bp:p2:triangular}(d)), hence so is $\Delta K=P\xi$.
\end{proof}

\begin{lemma}[The new invariance error]\label{bp:p2:newton-error}
In the setting of Lemma~\ref{bp:p2:correction}, assume moreover
\[
\frac{C_{\Delta K}\varepsilon}{\gamma^2\delta^a}<\beta\qquad\text{and}\qquad
\frac{C_{\mathrm{sym}}\varepsilon}{\gamma\delta^{\varkappa+1}}\le\frac12.
\]
Define $\bar K:=K+\Delta K$ on $T_{r-2\delta}$ and $\bar K:=K$ elsewhere. Then $\bar
K(T_{r-2\delta})\subset\mathfrak B$; the matrix $S:=I_{2n}-\Omega_0E_{\mathrm{sym}}$ is invertible with
$\abs{S^{-1}}\le2$ on $T_{r-2\delta}$; and on $T_{r-2\delta}$
\[
\bar E=PS^{-1}E_{\mathrm{lin}}+\Delta^2X,\qquad E_{\mathrm{lin}}:=E_{\mathrm{red}}\xi+\Omega_0
E_{\mathrm{sym}}\partial_\omega\xi,\qquad\Delta^2X:=X_h(\bar K)-X_h(K)-A\Delta K.
\]
Consequently $\abs{\bar E}\le C_E\varepsilon^2/(\gamma^4\delta^{2a})$ on $T_{r-2\delta}$.
\source{HL19 (5.1)--(5.4), p.~31; (5.29)--(5.32), pp.~34--35. Cross-check: GJLV05 Prop.~22, p.~893.}
\sourcegap{HL19 (5.30) inverts $-\Omega_0P^\top\Omega(K)$ without stating why it is invertible; the
argument is supplied in the first step below.}
\end{lemma}

\begin{proof}
\proofstep{Invertibility of $S$ and of $Q$}\label{step:p2:ne-invert}
On $T_{r-2\delta}$, $\abs{\Omega_0E_{\mathrm{sym}}}\le\frac12$ (Lemma~\ref{bp:p2:geometry}), so
$S$ is invertible with $\abs{S^{-1}}\le2$ by Lemma~\ref{bp:p1:neumann} ($A=I_{2n}$). Since $QP=S$
(Lemma~\ref{bp:p2:frame-identities}(f)), $Q(PS^{-1})=I_{2n}$; as $Q$ is square, $Q$ is invertible and
$Q^{-1}=PS^{-1}$.

\proofstep{The corrected torus stays in the domain}\label{step:p2:ne-domain}
For $\theta\in T_{r-2\delta}$, $\abs{\Delta K(\theta)}<\beta$, so $\bar K(\theta)\in\mathbb
B_\beta(K(\theta))\subset\mathfrak B$ by \ref{bp:p2:A2}, and the segment from $K(\theta)$ to
$\bar K(\theta)$ lies in this ball.

\proofstep{Expansion of the new error}\label{step:p2:ne-expand}
On $T_{r-2\delta}$, $\bar K$ is holomorphic and
$\bar E=X_h(K+\Delta K)-L\omega-\partial_\omega\Delta K=E+A\Delta K-\partial_\omega\Delta K+\Delta^2X$.
By the product rule, $A\Delta K-\partial_\omega\Delta K=(AP-\partial_\omega P)\xi-P\partial_\omega\xi$.

\proofstep{Reduction by $Q$}\label{step:p2:ne-reduce}
By Lemma~\ref{bp:p2:frame-identities}(e), (f),
$Q(E+A\Delta K-\partial_\omega\Delta K)=-\eta+(\Lambda+E_{\mathrm{red}})\xi-(I_{2n}-\Omega_0
E_{\mathrm{sym}})\partial_\omega\xi=(\Lambda\xi-\partial_\omega\xi-\eta)+E_{\mathrm{lin}}=E_{\mathrm{lin}}$,
using Lemma~\ref{bp:p2:triangular}(b). Applying $Q^{-1}=PS^{-1}$ gives the formula for $\bar E$.

\proofstep{Estimate of the linear part}\label{step:p2:ne-lin}
On $T_{r-2\delta}$, by Lemmas~\ref{bp:p2:geometry} and~\ref{bp:p2:correction},
$\abs{E_{\mathrm{red}}\xi}\le C_{\mathrm{red}}C_\xi\varepsilon^2/(\gamma^2\delta^{a+1})$ and
$\abs{\Omega_0E_{\mathrm{sym}}\partial_\omega\xi}\le C_{\mathrm{sym}}C_g\varepsilon^2/(\gamma^2
\delta^{2\varkappa+3})$; since $2\varkappa+3=a+1$, $\abs{E_{\mathrm{lin}}}\le C_{\mathrm{lin}}
\varepsilon^2/(\gamma^2\delta^{a+1})$.

\proofstep{Estimate of the quadratic part}\label{step:p2:ne-quad}
By Lemma~\ref{bp:p2:hamiltonian}(e) with $z=K(\theta)$, $w=\Delta K(\theta)$:
$\abs{\Delta^2X}\le M\abs{\Delta K}^2\le\mathfrak mC_{\Delta K}^2\varepsilon^2/(\gamma^4\delta^{2a})$.

\proofstep{The quadratic estimate}\label{step:p2:ne-total}
With $\abs P\le2\mathfrak m^2$ and $\abs{S^{-1}}\le2$,
\[
\abs{\bar E}\le\abs P\abs{S^{-1}}\abs{E_{\mathrm{lin}}}+\abs{\Delta^2X}\le
\frac{4\mathfrak m^2C_{\mathrm{lin}}\varepsilon^2}{\gamma^2\delta^{a+1}}+
\frac{\mathfrak mC_{\Delta K}^2\varepsilon^2}{\gamma^4\delta^{2a}},
\]
and $\gamma^2\delta^{a-1}\le\Theta^2$ (Lemma~\ref{bp:p2:frequency}, $\delta\le1$) turns the first term into
$4\mathfrak m^2\Theta^2C_{\mathrm{lin}}\varepsilon^2/(\gamma^4\delta^{2a})$.
\end{proof}

%----------------------------------------------------------------------------------------------------
\subsection{Perturbation of the frame and the torsion}\label{bp:p2:sec-perturbation}

The correction changes $K$ by $d=O(\varepsilon/(\gamma^2\delta^a))$. The nondegeneracy data change
accordingly: $DK$, the Gram inverse and the normal frame by $O(d/\delta)$, the torsion by
$O(d/\delta^2)$. The next lemma quantifies this for any holomorphic perturbation; it does not use the
invariance error, and it is also the stability statement needed in Part~\ref{bp:p4} (node STAB of
\texttt{docs/PLAN.md}).

\begin{lemma}[Perturbation of the frame and the torsion]\label{bp:p2:perturbation}
Let $r>0$ and $0<\delta\le1$ with $2\delta<r$; let $\abs\omega\le\Theta$. Let $K$ be holomorphic on
$T_r$ with $\mathbb B_\beta(K(\theta))\subset\mathfrak B$ for $\theta\in T_r$, $\abs{DK}\le\sigma_K$,
$\det G\ne0$ and $\abs B\le\sigma_B$ on $T_r$, and $\avg T$ invertible with $\abs{\avg T^{-1}}\le
\sigma_T$. Let $\tilde K$ be holomorphic on $T_r$ with $\abs{\tilde K-K}\le d$ on $T_r$, where
\[
d<\beta\qquad\text{and}\qquad2\mathfrak mC_{\Delta T}\,d/\delta^2\le1.
\]
Then:
\begin{enumerate}[label=(\alph*)]
\item $\mathbb B_{\beta-d}(\tilde K(\theta))\subset\mathfrak B$ for $\theta\in T_r$;
\item $\abs{D\tilde K-DK}\le nd/\delta$ on $T_{r-\delta}$;
\item $\det\tilde G\ne0$, $\abs{\tilde B}\le2\mathfrak m$ and $\abs{\tilde B-B}\le C_Bd/\delta$ on
  $T_{r-\delta}$;
\item $\abs{\tilde N-N}\le C_{\Delta N}d/\delta$ and $\abs{\tilde N}\le2\mathfrak m^2$ on $T_{r-\delta}$;
\item $\abs{\tilde T-T}\le C_{\Delta T}d/\delta^2$ on $T_{r-2\delta}$;
\item $\tilde T$ is continuous on $\R^n$, $\avg{\tilde T}$ is invertible, and $\abs{\avg{\tilde T}^{-1}
  -\avg T^{-1}}\le2\mathfrak m^2C_{\Delta T}d/\delta^2$.
\end{enumerate}
\source{HL19 Lemma~5.2, Steps 3--7, (5.34)--(5.67), pp.~36--42; Lemma~A.1, p.~55.}
\sourcegap{HL19 Step~7 (p.~41) skips the intermediate computations for the new torsion; they are
replaced here by Cauchy estimates and supplied in full.}
\end{lemma}

\begin{proof}
\proofstep{Consequences of the smallness}\label{step:p2:pt-small}
From the definitions, $C_{\Delta T}\ge2nC_{\Delta N}C_X\ge10n^3C_{\Delta N}$, $C_{\Delta N}\ge\max\{2n
\mathfrak m,\mathfrak mC_B\}$ and $\mathfrak mC_B=2\mathfrak m^3C_G\ge C_G$; hence $C_{\Delta T}\ge
\max\{n,C_G,C_{\Delta N}\}$. Since $\delta\le1$, $d/\delta\le d/\delta^2\le1/(2\mathfrak mC_{\Delta T})$, and
therefore
\[
\frac{nd}\delta\le\frac1{2\mathfrak m}\le\mathfrak m,\qquad\frac{\mathfrak mC_Gd}\delta\le\frac12,\qquad
\frac{C_{\Delta N}d}\delta\le\frac1{2\mathfrak m}\le\mathfrak m^2,\qquad
\frac{\mathfrak mC_{\Delta T}d}{\delta^2}\le\frac12.
\]

\proofstep{Domain}\label{step:p2:pt-domain}
If $\abs{z-\tilde K(\theta)}<\beta-d$ then $\abs{z-K(\theta)}<\beta$, so $z\in\mathfrak B$. In
particular the segment $[K(\theta),\tilde K(\theta)]$ lies in $\mathbb B_\beta(K(\theta))$, and
Lemma~\ref{bp:p2:hamiltonian}(e) gives $\abs{\tilde A-A}\le Md\le\mathfrak md$ on $T_r$.

\proofstep{Tangent frame}\label{step:p2:pt-dl}
Lemma~\ref{bp:p2:cauchy}(a) applied to $\tilde K-K$ gives $\abs{\tilde L-L}\le nd/\delta$ and
$\abs{(\tilde L-L)^\top}\le2nd/\delta$ on $T_{r-\delta}$. Hence $\abs{\tilde L}\le2\mathfrak m$ and
$\abs{\tilde L^\top}\le2n\mathfrak m+2nd/\delta\le4n\mathfrak m$.

\proofstep{Gram matrix}\label{step:p2:pt-gram}
$\tilde G-G=\tilde L^\top(\tilde L-L)+(\tilde L-L)^\top L$, so $\abs{\tilde G-G}\le4n\mathfrak m\cdot
nd/\delta+2nd/\delta\cdot\mathfrak m\le C_Gd/\delta$ on $T_{r-\delta}$.

\proofstep{Gram inverse (invertibility before use)}\label{step:p2:pt-binv}
On $T_{r-\delta}$, $\abs B\abs{\tilde G-G}\le\mathfrak mC_Gd/\delta\le\frac12$. By
Lemma~\ref{bp:p1:neumann}, $\tilde G$ is invertible, $\abs{\tilde B}\le2\abs B\le2\mathfrak m$ and
$\abs{\tilde B-B}\le2\mathfrak m^2C_Gd/\delta=C_Bd/\delta$. This is (c).

\proofstep{Normal frame}\label{step:p2:pt-n}
$\tilde N-N=\Omega_0\bigl((\tilde L-L)\tilde B+L(\tilde B-B)\bigr)$, so $\abs{\tilde N-N}\le
(2n\mathfrak m+\mathfrak mC_B)d/\delta=C_{\Delta N}d/\delta$ and $\abs{\tilde N}\le\mathfrak m^2+
C_{\Delta N}d/\delta\le2\mathfrak m^2$ on $T_{r-\delta}$. In particular $\abs{\tilde N^\top}\le
4n\mathfrak m^2$ and $\abs{(\tilde N-N)^\top}\le2nC_{\Delta N}d/\delta$.

\proofstep{Displacement of the normal frame}\label{step:p2:pt-xn}
On $T_{r-2\delta}$: $\tilde X_{\tilde N}-X_N=(\tilde A-A)\tilde N+A(\tilde N-N)-\partial_\omega(\tilde
N-N)$. The first two terms are bounded by $2\mathfrak m^3d+\mathfrak mC_{\Delta N}d/\delta$; the third,
by Lemma~\ref{bp:p2:cauchy}(b) on $T_{r-\delta}$, by $n\Theta C_{\Delta N}d/\delta^2$. With $\delta\le1$
the sum is $\le C_{\Delta X}d/\delta^2$.

\proofstep{Torsion}\label{step:p2:pt-torsion}
$\tilde T-T=\tilde N^\top\Omega_0(\tilde X_{\tilde N}-X_N)+(\tilde N-N)^\top\Omega_0X_N$. With
$\abs{X_N}\le C_X/\delta$ on $T_{r-\delta}$ (Lemma~\ref{bp:p2:frame-basic}(d), whose hypotheses are those
of the present lemma), on $T_{r-2\delta}$:
$\abs{\tilde T-T}\le4n\mathfrak m^2C_{\Delta X}d/\delta^2+2nC_{\Delta N}C_Xd/\delta^2=C_{\Delta T}d/
\delta^2$.

\proofstep{Twist (integrability and invertibility before use)}\label{step:p2:pt-twist}
By (c), $\tilde N$ and hence $\tilde T$ are holomorphic on $T_{r-\delta}\supset\R^n$, so $\tilde T$ is
continuous on $\R^n$ and $\avg{\tilde T}$ is defined. By Lemma~\ref{bp:p2:averages}(a),
$\abs{\avg{\tilde T}-\avg T}\le C_{\Delta T}d/\delta^2$, and $\abs{\avg T^{-1}}\abs{\avg{\tilde T}-\avg
T}\le\mathfrak mC_{\Delta T}d/\delta^2\le\frac12$. Lemma~\ref{bp:p1:neumann} gives (f).
\end{proof}

%----------------------------------------------------------------------------------------------------
\subsection{The KAM step}\label{bp:p2:sec-step}

\begin{proposition}[The KAM step]\label{bp:p2:step}
Let $r>0$, $0<\delta\le1$ with $4\delta<r$, $\beta>0$ and $0\le\varepsilon\le1$. Let $K$ be
$(r,\beta,\varepsilon)$-admissible, and assume
\[
\mathfrak c_{\mathrm{step}}\,\frac{\varepsilon}{\gamma^4\delta^{2a}}\le\min\{1,\beta\}.
\]
Let $\bar K$ be the map of Lemma~\ref{bp:p2:newton-error} and $d:=C_{\Delta K}\varepsilon/(\gamma^2
\delta^a)$. Then:
\begin{enumerate}[label=(\alph*)]
\item $\bar K$ is holomorphic and primary on $T_{r-2\delta}$ and real on reals;
\item $\abs{\bar K-K}\le d\le\beta/2$ on $T_{r-2\delta}$, and $\mathbb B_{\beta-d}(\bar K(\theta))\subset
  \mathfrak B$ for $\theta\in T_{r-2\delta}$;
\item $\abs{\bar E}\le C_E\varepsilon^2/(\gamma^4\delta^{2a})$ on $T_{r-2\delta}$;
\item $\abs{D\bar K-DK}\le nd/\delta$ on $T_{r-3\delta}$;
\item $\det\bar G\ne0$ and $\abs{\bar B-B}\le C_Bd/\delta$ on $T_{r-3\delta}$;
\item $\avg{\bar T}$ is invertible and $\abs{\avg{\bar T}^{-1}-\avg T^{-1}}\le2\mathfrak m^2C_{\Delta T}
  d/\delta^2$.
\end{enumerate}
\source{HL19 Lemma~5.2, (5.8)--(5.18), pp.~32--33, with (5.24), (5.25), (5.32), (5.34), (5.41),
(5.67). Cross-check: GJLV05 \S8 and the iterative step of \S5.}
\end{proposition}

\begin{proof}
\proofstep{The smallness conditions used}\label{step:p2:st-small}
By Lemma~\ref{bp:p2:frequency}, $\gamma\le\Theta$; with $\delta\le1$, $\gamma^i\delta^j\le\Theta^i$ for
$i,j\ge0$. Hence, since $\varkappa+1\le2a$ and $a+2\le2a$:
$C_{\mathrm{sym}}\varepsilon/(\gamma\delta^{\varkappa+1})\le\Theta^3C_{\mathrm{sym}}\varepsilon/(\gamma^4
\delta^{2a})\le\frac12$; $d\le\Theta^2C_{\Delta K}\varepsilon/(\gamma^4\delta^{2a})\le\beta/2$; and
$2\mathfrak mC_{\Delta T}d/\delta^2\le2\mathfrak m\Theta^2C_{\Delta T}C_{\Delta K}\varepsilon/(\gamma^4
\delta^{2a})\le\frac12$. Each follows from $\mathfrak c_{\mathrm{step}}\varepsilon/(\gamma^4\delta^{2a})\le
\min\{1,\beta\}$ and $\Theta\ge1$.

\proofstep{Correction and new error}\label{step:p2:st-error}
Lemma~\ref{bp:p2:correction} applies (as $2\delta<r$), and Lemma~\ref{bp:p2:newton-error} applies by the
previous step, since $d\le\beta/2<\beta$. They give (a) (with Lemma~\ref{bp:p2:correction} for
periodicity and realness of $\Delta K$, so $\bar K(\theta+k)=K(\theta)+(k,0)+\Delta K(\theta)$), the
bound in (b), and (c).

\proofstep{Ball margin}\label{step:p2:st-ball}
If $\abs{z-\bar K(\theta)}<\beta-d$, then $\abs{z-K(\theta)}<\beta$, so $z\in\mathfrak B$ by \ref{bp:p2:A2}.

\proofstep{Perturbation on the strip $T_{r-2\delta}$}\label{step:p2:st-perturb}
On $T_{r-2\delta}$, $K$ satisfies the hypotheses of Lemma~\ref{bp:p2:perturbation} (restriction of
\ref{bp:p2:A2}--\ref{bp:p2:A5}; $\abs\omega\le\Theta$ by Lemma~\ref{bp:p2:frequency}), and $\tilde K:=\bar
K$ is holomorphic there with $\abs{\bar K-K}\le d<\beta$ and $2\mathfrak mC_{\Delta T}d/\delta^2\le1$ by
the first step. Since $2\delta<r-2\delta$, the lemma gives (d) and (e) on $T_{r-3\delta}$ and (f).
\end{proof}

%----------------------------------------------------------------------------------------------------
\subsection{The iteration}\label{bp:p2:sec-iteration}

\begin{definition}[Schedule of the iteration]\label{bp:p2:schedule}
Let $0<\rho_\infty<\rho$, $\gamma>0$ and $e\ge0$. With $\delta_0$ from
Definition~\ref{bp:p2:constants} put, for $s\ge0$,
\[
\delta_s:=2^{-s}\delta_0,\qquad r_s:=\rho-8\delta_0(1-2^{-s}),\qquad\varepsilon_0:=e,\quad
\varepsilon_{s+1}:=\frac{C_E\,\varepsilon_s^2}{\gamma^4\delta_s^{2a}},\qquad
d_s:=\frac{C_{\Delta K}\varepsilon_s}{\gamma^2\delta_s^a}.
\]
Then $r_0=\rho$, $r_{s+1}=r_s-4\delta_s$, $\rho_\infty\le\rho-8\delta_0<r_s$, $\delta_s\le1$, and
$4\delta_s<r_s$ for every $s$.
\source{HL19 \S5.3, p.~42 (there $\delta_s=\delta_0a_1^{-s}$ and $\rho_s=\rho_{s-1}-3\delta_{s-1}$; here
$a_1=2$ and the loss per step is $4\delta_s$).}
\end{definition}

The closeness bound must be linear in $e$, so we need the decay of the errors in a form proportional to
the initial error; part (a) below is a sharper form of Lemma~\ref{bp:p1:iteration}.

\begin{lemma}[Decay of the errors]\label{bp:p2:error-decay}
\begin{enumerate}[label=(\alph*)]
\item Let $Q>0$, $p\ge0$ and $y_0,y_1,\dots\ge0$ with $y_{m+1}\le Q2^{pm}y_m^2$ for every $m\ge0$. If
  $q:=Q2^py_0\le1$, then $y_m\le q^{2^m-1}2^{-pm}y_0$ for every $m\ge0$. If moreover $q\le\frac12$, then
  $y_m\le Q^{-1}2^{-p(m+1)}2^{-2^m}$, which is the conclusion of Lemma~\ref{bp:p1:iteration}.
\item If $q:=2^{2a}C_E\,e/(\gamma^4\delta_0^{2a})\le1$, then $\varepsilon_s\le q^{2^s-1}2^{-2as}e\le
  2^{-2as}e$ for all $s\ge0$. Consequently $\varepsilon_s/\delta_s^j\le2^{-(2a-j)s}e/\delta_0^j\le
  e/\delta_0^j$ for $0\le j\le2a$, and $\sum_{s\ge0}\varepsilon_s/\delta_s^j\le2e/\delta_0^j$ for
  $0\le j\le2a-1$.
\end{enumerate}
\source{HL19 (5.68)--(5.69), p.~43; Lemma~\ref{bp:p1:iteration}.}
\end{lemma}

\begin{proof}
\proofstep{Induction}\label{step:p2:ed-induction}
The case $m=0$ is trivial. If $y_m\le q^{2^m-1}2^{-pm}y_0$, then
$y_{m+1}\le Q2^{pm}q^{2^{m+1}-2}2^{-2pm}y_0^2=(Q2^py_0)\,q^{2^{m+1}-2}2^{-p(m+1)}y_0=
q^{2^{m+1}-1}2^{-p(m+1)}y_0$.

\proofstep{Comparison with Part~1}\label{step:p2:ed-part1}
If $0<q\le\frac12$, then $q^{2^m-1}2^{-pm}y_0=q^{2^m}2^{-pm}y_0/q=Q^{-1}2^{-p(m+1)}q^{2^m}\le
Q^{-1}2^{-p(m+1)}2^{-2^m}$; if $q=0$, all $y_m$ vanish.

\proofstep{The errors of the schedule}\label{step:p2:ed-errors}
By Definition~\ref{bp:p2:schedule} and $\delta_s^{-2a}=2^{2as}\delta_0^{-2a}$, the errors satisfy
$\varepsilon_{s+1}=Q2^{2as}\varepsilon_s^2$ with $Q=C_E/(\gamma^4\delta_0^{2a})$. Part (a) with $p=2a$ and
$y_0=e$ gives $\varepsilon_s\le q^{2^s-1}2^{-2as}e$, and $q\le1$ gives the second inequality.

\proofstep{Sums}\label{step:p2:ed-sums}
$\varepsilon_s\delta_s^{-j}\le2^{-2as}e\,2^{js}\delta_0^{-j}$; for $2a-j\ge1$ the geometric series
$\sum_s2^{-(2a-j)s}$ is at most $2$.
\end{proof}

\begin{proposition}[The induction]\label{bp:p2:iteration}
Let $0<\rho_\infty<\rho$, $\mu>0$, and let $K$ satisfy, on $T_\rho$, the hypotheses H2--H4 of
Theorem~\ref{bp:p0:g1}: $K$ holomorphic, primary and real on reals; $\mathbb B_\mu(K(\theta))\subset
\mathfrak B$; $\abs{DK}\le\sigma_K-\mu$; $\det G\ne0$ and $\abs B\le\sigma_B-\mu$; $\avg T$ invertible with
$\abs{\avg T^{-1}}\le\sigma_T-\mu$. Let $\abs E\le e$ on $T_\rho$ with
\[
e\le1\qquad\text{and}\qquad\mathfrak c_*\,\frac{e}{\gamma^4\delta_0^{2a}}\le1.
\]
Then there are maps $K_s:\C^n\to\C^n\times\C^n$, $K_0=K$, such that, with $D_s:=\sum_{j<s}d_j$, for every
$s\ge0$:
\begin{enumerate}[label=(\roman*)]
\item $K_s$ is $(r_s,\mu-D_s,\varepsilon_s)$-admissible, and $\varepsilon_s\le e\le1$;
\item $D_s\le\mu/4$, and $\abs{K_s-K}\le D_s$ on $T_{r_s}$;
\item $\abs{DK_s}\le\sigma_K-\mu+\sum_{j<s}nd_j/\delta_j\le\sigma_K-\mu/2$ on $T_{r_s}$;
\item $\abs{B_s}\le\sigma_B-\mu+\sum_{j<s}C_Bd_j/\delta_j\le\sigma_B-\mu/2$ on $T_{r_s}$;
\item $\abs{\avg{T_s}^{-1}}\le\sigma_T-\mu+\sum_{j<s}2\mathfrak m^2C_{\Delta T}d_j/\delta_j^2\le\sigma_T-
  \mu/2$;
\item $\abs{K_{s+1}-K_s}\le d_s$ and $\abs{DK_{s+1}-DK_s}\le nd_s/\delta_s$ on $T_{r_{s+1}}$.
\end{enumerate}
\source{HL19 \S5.3, (5.68)--(5.73), pp.~42--44.}
\end{proposition}

\begin{proof}
\proofstep{Summed bounds}\label{step:p2:it-sums}
By Lemma~\ref{bp:p2:frequency} applied to $K$ ($e\le1$), $\gamma\le\Theta$, so $e/(\gamma^2\delta_0^j)\le
\Theta^2e/(\gamma^4\delta_0^{2a})$ for $j\le2a$. Since $\mathfrak c_*\ge2^{2a}C_E$,
Lemma~\ref{bp:p2:error-decay} applies; with $j=a,a+1,a+2\le2a-1$ (as $a\ge3$) and
$\mathfrak c_{\mathrm{sum}}e/(\gamma^4\delta_0^{2a})\le\mu/2$ (as $\mathfrak c_*\ge2\mathfrak
c_{\mathrm{sum}}/\mu$):
\begin{align*}
\textstyle\sum_jd_j&\le2\Theta^2C_{\Delta K}e/(\gamma^4\delta_0^{2a})\le\mu/4,\\
\textstyle\sum_jnd_j/\delta_j&\le2n\Theta^2C_{\Delta K}e/(\gamma^4\delta_0^{2a})\le\mu/4,\\
\textstyle\sum_jC_Bd_j/\delta_j&\le2\Theta^2C_BC_{\Delta K}e/(\gamma^4\delta_0^{2a})\le\mu/4,\\
\textstyle\sum_j2\mathfrak m^2C_{\Delta T}d_j/\delta_j^2&\le4\Theta^2\mathfrak m^2C_{\Delta T}C_{\Delta K}
e/(\gamma^4\delta_0^{2a})\le\mu/2.
\end{align*}
These give the right-hand inequalities in (ii)--(v) for every $s$.

\proofstep{Step smallness at every stage}\label{step:p2:it-stepsmall}
By Lemma~\ref{bp:p2:error-decay}, $\varepsilon_s\le e\le1$ and $\varepsilon_s/\delta_s^{2a}\le e/
\delta_0^{2a}$, so $\mathfrak c_{\mathrm{step}}\varepsilon_s/(\gamma^4\delta_s^{2a})\le\mathfrak
c_{\mathrm{step}}e/(\gamma^4\delta_0^{2a})\le\min\{1,\mu/2\}$, as $\mathfrak c_*\ge\mathfrak
c_{\mathrm{step}}\max\{1,2/\mu\}$. Since $\mu-D_s\ge\mu/2$ whenever (ii) holds, the smallness hypothesis of
Proposition~\ref{bp:p2:step} holds with $(r,\delta,\beta,\varepsilon)=(r_s,\delta_s,\mu-D_s,
\varepsilon_s)$.

\proofstep{Base case}\label{step:p2:it-base}
$K_0=K$ is $(\rho,\mu,e)$-admissible since $\sigma_K-\mu\le\sigma_K$ etc., and (ii)--(v) hold with empty
sums.

\proofstep{Inductive step}\label{step:p2:it-step}
Assume (i)--(v) at $s$. Proposition~\ref{bp:p2:step} applies ($4\delta_s<r_s$, $\delta_s\le1$) and
defines $K_{s+1}:=\bar K_s$. On $T_{r_{s+1}}\subset T_{r_s-3\delta_s}\subset T_{r_s-2\delta_s}$ it gives:
\ref{bp:p2:A1}; the ball $\mathbb B_{\mu-D_s-d_s}(K_{s+1}(\theta))=\mathbb B_{\mu-D_{s+1}}(K_{s+1}
(\theta))\subset\mathfrak B$; $\abs{E_{s+1}}\le\varepsilon_{s+1}$; and (vi). Then
$\abs{K_{s+1}-K}\le D_s+d_s=D_{s+1}$; $\abs{DK_{s+1}}\le\abs{DK_s}+nd_s/\delta_s$; $\det G_{s+1}\ne0$
and $\abs{B_{s+1}}\le\abs{B_s}+C_Bd_s/\delta_s$; $\avg{T_{s+1}}$ invertible with $\abs{\avg{T_{s+1}}^{-1}}
\le\abs{\avg{T_s}^{-1}}+2\mathfrak m^2C_{\Delta T}d_s/\delta_s^2$. With the summed bounds this is
(ii)--(v) at $s+1$, and these bounds give \ref{bp:p2:A3}--\ref{bp:p2:A5}, so (i) holds at $s+1$.
\end{proof}

\begin{proposition}[The limit torus]\label{bp:p2:limit}
In the setting of Proposition~\ref{bp:p2:iteration}, define
\[
K_\infty(\theta):=\lim_{s\to\infty}K_s(\theta)\quad(\theta\in T_{\rho_\infty}),\qquad
K_\infty(\theta):=K(\theta)\quad(\theta\notin T_{\rho_\infty}).
\]
The limit exists by (a) below. Then:
\begin{enumerate}[label=(\alph*)]
\item $K_s\to K_\infty$ and $DK_s\to DK_\infty$ uniformly on $T_{\rho_\infty}$, and $K_\infty$ is
  holomorphic on $T_{\rho_\infty}$;
\item $K_\infty$ is primary on $T_{\rho_\infty}$ and real on reals;
\item $\abs{K_\infty-K}\le2C_{\Delta K}e/(\gamma^2\delta_0^a)\le\mu/4$ on $T_{\rho_\infty}$;
\item $\mathbb B_{\mu/2}(K_\infty(\theta))\subset\mathfrak B$ and $\abs{DK_\infty(\theta)}\le\sigma_K-\mu/2$
  for $\theta\in T_{\rho_\infty}$;
\item $E_\infty(\theta)=X_h(K_\infty(\theta))-DK_\infty(\theta)\omega=0$ for $\theta\in T_{\rho_\infty}$.
\end{enumerate}
\source{HL19 p.~44 and (5.78), p.~44. Cross-check: GJLV05 Theorem~5, p.~865.}
\sourcegap{HL19 (p.~44) state that ``the iterative scheme converges to a true quasi-periodic torus
$K_\infty$'' with $K_\infty\in A(\T^d_{\rho_\infty})$ and the bounds on $DK_\infty$, without the
convergence argument; it is supplied here.}
\end{proposition}

\begin{proof}
\proofstep{Uniform convergence of $K_s$}\label{step:p2:lim-k}
$T_{\rho_\infty}\subset T_{r_s}$ for every $s$ (Definition~\ref{bp:p2:schedule}). By
Proposition~\ref{bp:p2:iteration}(vi), $\abs{K_{s+1}-K_s}\le d_s$ on $T_{\rho_\infty}$, and
$\sum_sd_s\le2C_{\Delta K}e/(\gamma^2\delta_0^a)<\infty$ (Lemma~\ref{bp:p2:error-decay} with $j=a$).
Hence $(K_s)$ is uniformly Cauchy on $T_{\rho_\infty}$ and converges uniformly to $K_\infty$.

\proofstep{Uniform convergence of $DK_s$}\label{step:p2:lim-dk}
Likewise $\abs{DK_{s+1}-DK_s}\le nd_s/\delta_s$ on $T_{\rho_\infty}$ with $\sum_snd_s/\delta_s<\infty$, so
$DK_s$ converges uniformly on $T_{\rho_\infty}$ to some $\Lambda_\infty$ with values in $\C$-linear maps.

\proofstep{Holomorphy of the limit}\label{step:p2:lim-holo}
Each $K_s$ is holomorphic on the open set $T_{\rho_\infty}$ with derivative $DK_s$. A sequence of
differentiable maps on an open set whose derivatives converge uniformly and which converges pointwise has
a differentiable limit whose derivative is the limit of the derivatives (Mathlib
\Lean{hasFDerivAt\_of\_tendstoUniformlyOn}). Hence $K_\infty$ is complex differentiable on
$T_{\rho_\infty}$ with $DK_\infty=\Lambda_\infty$; this is (a).

\proofstep{Primary and real}\label{step:p2:lim-primary}
For $\theta\in T_{\rho_\infty}$ and $k\in\Z^n$, $\theta+k\in T_{\rho_\infty}$ and $K_s(\theta+k)=K_s(\theta)
+(k,0)$ for all $s$; let $s\to\infty$. For $x\in\R^n\subset T_{\rho_\infty}$, $K_s(x)\in\R^n\times\R^n$, a
closed set, so $K_\infty(x)$ is real.

\proofstep{Closeness}\label{step:p2:lim-close}
On $T_{\rho_\infty}$, $\abs{K_\infty-K}=\lim_s\abs{K_s-K}\le\sup_sD_s\le\sum_sd_s$. By the summed bounds
in the proof of Proposition~\ref{bp:p2:iteration}, $\sum_sd_s\le2C_{\Delta K}e/(\gamma^2\delta_0^a)\le
\mu/4$.

\proofstep{Margins}\label{step:p2:lim-margins}
If $\abs{z-K_\infty(\theta)}<\mu/2$, then $\abs{z-K(\theta)}<\mu/2+\mu/4<\mu$, so $z\in\mathfrak B$ by
hypothesis H2. And $\abs{DK_\infty(\theta)}=\lim_s\abs{DK_s(\theta)}\le\sigma_K-\mu/2$ by
Proposition~\ref{bp:p2:iteration}(iii).

\proofstep{Invariance}\label{step:p2:lim-invariance}
Fix $\theta\in T_{\rho_\infty}$. Then $K_\infty(\theta)\in\mathfrak B$, where $X_h$ is continuous, so
$E_s(\theta)=X_h(K_s(\theta))-DK_s(\theta)\omega\to X_h(K_\infty(\theta))-DK_\infty(\theta)\omega$. Since
$\abs{E_s(\theta)}\le\varepsilon_s\le2^{-2as}e\to0$ (Proposition~\ref{bp:p2:iteration}(i) and
Lemma~\ref{bp:p2:error-decay}), $E_\infty(\theta)=0$.
\end{proof}

%----------------------------------------------------------------------------------------------------
\subsection{Proof of Theorem G1}\label{bp:p2:sec-proof}

\begin{theorem}[G1 with explicit constants]\label{bp:p2:g1}
Let $n\ge2$, $\tau\ge n-1$, $0<\rho_\infty<\rho$, and let $M$, $\sigma_K$, $\sigma_B$, $\sigma_T$, $\mu$,
$\bar\gamma$ be positive. With the notation of Definition~\ref{bp:p2:constants} put
\[
C_1:=\rho^{4\tau}\max\bigl\{\bar\gamma^4,\ \mathfrak c_*\delta_0^{-2a}\bigr\},\qquad
C_2:=2C_{\Delta K}\,\rho^{2\tau}\delta_0^{-a}.
\]
Then $C_1,C_2>0$, and the conclusion of Theorem~\ref{bp:p0:g1} holds with these constants for all
$\mathfrak B$, $h$, $K$, $\omega$, $\gamma$, $e$ satisfying its hypotheses. Here $C_1$ depends only on
$(n,\tau,\rho,\rho_\infty,M,\sigma_K,\sigma_B,\sigma_T,\mu,\bar\gamma)$, and $C_2$ only on
$(n,\tau,\rho,\rho_\infty,M,\sigma_K,\sigma_B,\sigma_T)$.
\frozen{KAM.aposteriori\_kam}
\source{HL19 Theorem~3.1, pp.~16--18; (5.74)--(5.78), p.~44 (ledger D-G1-1, E-HL19-1).}
\end{theorem}

\begin{proof}
\proofstep{The constants}\label{step:p2:g1-constants}
All constants of Definition~\ref{bp:p2:constants} are finite and $\ge1$, $\delta_0\in(0,1]$ and
$\rho>0$, so $C_1\ge\rho^{4\tau}\bar\gamma^4>0$ and $C_2\ge2\rho^{2\tau}>0$. The dependence is read off
from the definitions; the constant $c(n,\tau)$ of Lemma~\ref{bp:p1:cohomological} enters through $c_R$.

\proofstep{The standing assumptions}\label{step:p2:g1-standing}
Hypothesis H1 of Theorem~\ref{bp:p0:g1} (openness and translation invariance of $\mathfrak B$,
holomorphy, periodicity and realness of $h$, the bounds on $X_h$, $DX_h$, $D^2X_h$) and H5
($\gamma>0$, $\omega\in\Dio_{\gamma,\tau}$) are Assumption~\ref{bp:p2:standing}; the norms of the frozen
statement are the operator norms of Part~\ref{bp:p0} (Remark~\ref{bp:p2:frozen-agreement}).

\proofstep{Size of the error}\label{step:p2:g1-e}
$T_\rho\ne\emptyset$, so $0\le\abs{E(0)}\le e$. From $C_1e/(\gamma^4\rho^{4\tau})<1$ and $C_1\ge
\rho^{4\tau}\bar\gamma^4$ we get $e<\gamma^4/\bar\gamma^4\le1$, using $\gamma\le\bar\gamma$; and from
$C_1\ge\rho^{4\tau}\mathfrak c_*\delta_0^{-2a}$ we get $\mathfrak c_*e/(\gamma^4\delta_0^{2a})<1$.

\proofstep{Iteration and limit}\label{step:p2:g1-apply}
Hypotheses H2--H4 of Theorem~\ref{bp:p0:g1} are those of Proposition~\ref{bp:p2:iteration} (the
invertibility of $\det G$ and of $\avg T$ is part of H3, H4). With the previous step,
Propositions~\ref{bp:p2:iteration} and~\ref{bp:p2:limit} apply and produce $K_\infty$.

\proofstep{Checking the conclusion}\label{step:p2:g1-check}
By Proposition~\ref{bp:p2:limit}: $K_\infty$ is complex differentiable on $T_{\rho_\infty}$ (a), primary on
$T_{\rho_\infty}$ and real on reals (b); its invariance error vanishes on $T_{\rho_\infty}$ (e);
$\mathbb B_{\mu/2}(K_\infty(\theta))\subset\mathfrak B$ and $\abs{DK_\infty(\theta)}\le\sigma_K-\mu/2$ on
$T_{\rho_\infty}$ (d); and on $T_{\rho_\infty}$, by (c),
\[
\abs{K_\infty-K}\le\frac{2C_{\Delta K}e}{\gamma^2\delta_0^a}=\frac{C_2\,e}{\gamma^2\rho^{2\tau}}.
\]
These are the seven clauses of the conclusion of Theorem~\ref{bp:p0:g1}; the derivative, invariance error
and ball in them are the Fr\'echet derivative, \Lean{invarianceError} and the open ball of the sup metric,
which on the open strip $T_{\rho_\infty}$ are the objects used here (Remark~\ref{bp:p2:frozen-agreement}).
\end{proof}

\begin{remark}[Constants and their counterparts in HL19]\label{bp:p2:constants-register}
Each constant of Definition~\ref{bp:p2:constants} is introduced in one result:
\begin{center}\small
\begin{tabular}{lll}
constant & bounds & introduced in (HL19 counterpart)\\\hline
$\Theta$ & $\abs\omega$, $\gamma$ & Lemma~\ref{bp:p2:frequency} (none)\\
$C_X$, $C_T$ & $X_N$, $T$ & Lemma~\ref{bp:p2:frame-basic} ($C_{X_N}$, $C_T$: (4.59)--(4.60))\\
$C_\Omega$, $C_{\mathrm{sym}}$ & $\Omega_K$, $E_{\mathrm{sym}}$ & Lemma~\ref{bp:p2:geometry} ($C_{\Omega_K}$,
  $C_{\mathrm{sym}}$: (4.15), (4.39))\\
$C_{\mathrm{red}}$ & $E_{\mathrm{red}}$ & Lemma~\ref{bp:p2:geometry} ($C_{\mathrm{red}}$: (4.73))\\
$C_{\xi_0},C_{\xi N},C_g,C_\xi$ & $\xi^N_0,\xi^N,g^L,\xi$ & Lemma~\ref{bp:p2:correction}
  ((5.20)--(5.23), (5.27))\\
$C_{\Delta K}$ & $\Delta K$ & Lemma~\ref{bp:p2:correction} ($C_{\Delta K}$: (5.24))\\
$C_{\mathrm{lin}}$, $C_E$ & $E_{\mathrm{lin}}$, $\bar E$ & Lemma~\ref{bp:p2:newton-error} ((5.29),
  (5.32))\\
$C_G,C_B,C_{\Delta N}$ & $\tilde G-G$, $\tilde B-B$, $\tilde N-N$ & Lemma~\ref{bp:p2:perturbation}
  ((5.39), (5.41), (5.49))\\
$C_{\Delta X},C_{\Delta T}$ & $\tilde X_{\tilde N}-X_N$, $\tilde T-T$ & Lemma~\ref{bp:p2:perturbation}
  ((5.64)--(5.65))\\
$\mathfrak c_{\mathrm{step}}$ & step smallness & Proposition~\ref{bp:p2:step} ($\hat C_\Delta$, $C_E$: (5.8)--(5.9))\\
$\mathfrak c_{\mathrm{sum}}$, $\mathfrak c_*$ & summed changes & Proposition~\ref{bp:p2:iteration}
  ((5.74)--(5.77))\\
$C_1$, $C_2$ & G1 & Theorem~\ref{bp:p2:g1} ((5.74), (5.78))
\end{tabular}
\end{center}
Unlike HL19's constants, none of these depends on $\gamma$, on the gaps $\sigma-\norm\cdot$ of the
initial torus (these are replaced by the uniform margin $\mu$), or on the initial torus otherwise.
\end{remark}

%----------------------------------------------------------------------------------------------------
\subsection{Ledger candidates and deviations}\label{bp:p2:sec-ledger}

\paragraph{Errata candidates in HL19 (arXiv v1).} In addition to E-HL19-1 and E-HL19-2:
\begin{enumerate}[label=(L\arabic*)]
\item\label{bp:p2:L1} \emph{Lemma~A.1, p.~55.} Printed conclusion: ``$\abs{\bar M^{-1}-M^{-1}}<
  2\sigma^2\abs{\bar M-M}$''. False for $\bar M=M$, which satisfies the hypothesis (A.1) and gives
  $0<0$. Correction: $\le$. In (A.3), the printed ``$\frac{\sigma^2\abs{\bar M-M}}{\sigma-\abs{M^{-1}}}<
  \frac12$'' follows from (A.1) only with $\le$. The other conclusion $\abs{\bar M^{-1}}<\sigma$ is true
  (for $\bar M=M$ it is the hypothesis $\abs{M^{-1}}<\sigma$). Kind: typo (strict versus non-strict
  inequality). Effect: none here (we use Lemma~\ref{bp:p1:neumann}).
\item\label{bp:p2:L2} \emph{Lemma~5.2, (5.15)--(5.18), p.~33, and (5.68), p.~43.} Printed, e.g.,
  ``$\norm{\bar K-K}_{\rho-2\delta}<\frac{C_{\Delta K}}{\gamma^2\delta^{2\tau}}\norm E_\rho$'' and
  ``$\norm{\bar E}_{\rho-2\delta}<\frac{C_E}{\gamma^4\delta^{4\tau}}\norm E_\rho^2$''. False when $E=0$
  (both sides vanish); the same defect as E-HL19-1, which inherits it. Correction: $\le$. Kind: typo.
\item\label{bp:p2:L3} \emph{p.~42, sentence after (5.67).} Printed: ``and so, we obtain the estimates
  (5.12) and (5.16) on the new object.'' The estimates just proved are those on the torsion, (5.13) and
  (5.17); the printed sentence repeats the one after (5.41). Kind: typo.
\item\label{bp:p2:L4} \emph{p.~44, after (5.77).} Printed: ``$\norm{E_s}_{\rho_s}<a_1^{-4\tau s}\norm
  E_\rho\longrightarrow0$ when $s\to0$''. Should read $s\to\infty$. Kind: typo.
\item\label{bp:p2:L5} \emph{p.~44, end of the proof of Theorem~3.1.} Printed: ``$\operatorname{dist}
  (K_\infty(\T^d),\partial\mathcal B)>0$'', whereas Theorem~3.1 (p.~17) asserts
  ``$\operatorname{dist}(K_\infty(\T^d_{\rho_\infty}),\partial\mathcal B)>0$''. The proof establishes the
  latter (via (5.73)); the printed line is the weaker statement. Kind: typo.
\end{enumerate}
Not errata, recorded for readers: the proof of HL19 Lemma~4.3 (p.~22) needs a primitive $a$ with $a\circ K$
periodic; for primary tori this holds for $\alpha_0=\sum_iz_{n+i}\,dz_i$ but not for the second action
form of Rem.~2.1 (p.~6), $\frac12\sum_i(z_{n+i}dz_i-z_idz_{n+i})$. Rem.~2.3 assumes $a$ periodic, so the
text is consistent; Lemma~\ref{bp:p2:exactness} uses $\alpha_0=\sum_iI_i\,d\theta_i$.

\paragraph{Deviations from HL19 introduced in this part.}
\begin{enumerate}[label=(V\arabic*)]
\item\label{bp:p2:V1} \emph{Specialization:} $d=n$, canonical structure, Case~III ($A=0$), primary tori
  in $\T^n\times\C^n$ coordinates (refines D-COORD-1). In this case the reducibility defect is exactly
  $E_{\mathrm{red}}$ of Lemma~\ref{bp:p2:frame-identities}(e), of size $O(e/\delta)$.
\item\label{bp:p2:V2} \emph{Small-divisor exponent} $\varkappa=\tau+n$ (decision D12): the new error is
  $O(e^2/(\gamma^4\delta^{2a}))$ with $a=2\tau+2n+2$ instead of $O(e^2/(\gamma^4\delta^{4\tau}))$, the
  closeness $O(e/(\gamma^2\delta^a))$ instead of $O(e/(\gamma^2\delta^{2\tau}))$. The powers of $\gamma$
  are HL19's.
\item\label{bp:p2:V3} \emph{Open strips and pointwise bounds} instead of sup norms on closed strips and
  the spaces $A(\T^d_\rho)$; the limit's derivative is obtained from uniform convergence of $DK_s$ on the
  open strip $T_{\rho_\infty}$, so no extra shrinking is needed.
\item\label{bp:p2:V4} \emph{Frequency bound and Cauchy estimates} (Lemma~\ref{bp:p2:frequency},
  Remark~\ref{bp:p2:torsion-route}) instead of the $L_\omega$ calculus of HL19 Lemma~4.8 and Lemma~5.2
  Step~6; the torsion is $O(1/\delta)$ and the torsion perturbation is estimated by Cauchy estimates
  (Lemma~\ref{bp:p2:perturbation}) instead of HL19 Step~7. The strip loss per step is $4\delta_s$ instead
  of $3\delta_{s}$.
\item\label{bp:p2:V5} \emph{$\gamma$-independent step constants:} $\gamma\le\abs\omega\le\Theta$, so the
  constants do not grow with $\gamma$; the bound $\bar\gamma$ (ledger D-G1-1) is used only to force
  $e\le1$, through $C_1\ge\rho^{4\tau}\bar\gamma^4$.
\item\label{bp:p2:V6} \emph{Crude constants} through one size constant $\mathfrak m$; transposed bounds
  from $\abs{M^\top}\le p\abs M$ (Remark~\ref{bp:p2:transposes}), which removes HL19's $\sigma_{K^\top}$
  and $c_{X_h^\top,1}$; HL19's parameter $\nu$ and condition (4.40) are replaced by $e\le1$, $\delta\le1$.
\item\label{bp:p2:V7} \emph{Schedule:} $a_1=2$, $\delta_0=\min\{1,(\rho-\rho_\infty)/8\}$; the gap conditions
  of HL19 (5.9), which involve $\sigma-\norm{\cdot}$ of the current torus, are replaced by the uniform
  margin~$\mu$, and the conclusion keeps the margins $\mu/2$ (D-G1-1).
\item\label{bp:p2:V8} \emph{Unused hypothesis:} the bound $\abs{Dh}\le M$ of G1's H1 is not used.
\item\label{bp:p2:V9} \emph{Realness and periodicity} of the correction are proved
  (Lemmas~\ref{bp:p2:periodicity}, \ref{bp:p2:triangular}(d), \ref{bp:p2:correction}); HL19 build them into
  the spaces $A(\T^d_\rho)$ without comment.
\end{enumerate}

%----------------------------------------------------------------------------------------------------
\subsection{Transfer requests to Part 1}\label{bp:p2:sec-transfer}

Part~\ref{bp:p1} already provides the bounded-periodic, Taylor, matrix-norm, holomorphic-inverse and
real-torus-average lemmas used here (Lemmas~\ref{bp:p1:periodic-bounded}, \ref{bp:p1:taylor},
\ref{bp:p1:matrix-norms}, \ref{bp:p1:holomorphic-inverse}, \ref{bp:p1:averages-real}). The following
results are stated and proved in this part but belong to the infrastructure of Part~\ref{bp:p1}; they are
general and likely needed by Parts~\ref{bp:p3} and~\ref{bp:p4}.
\begin{enumerate}[label=(T\arabic*)]
\item Lemma~\ref{bp:p2:matrix}(a), (c), (d): the operator-norm characterization
  $\abs M=\sup_{\abs v\le1}\abs{Mv}$, block matrices, and the facts on $\Omega_0$.
\item Lemma~\ref{bp:p2:cauchy}: Cauchy bounds for $Df$, $(Df)^\top$ and $\partial_\omega f$, corollaries of
  Lemma~\ref{bp:p1:cauchy-strip}.
\item Lemma~\ref{bp:p2:averages}(a): averages of continuous periodic matrix-valued functions on $\R^n$,
  with $\abs{\avg f}\le\sup_{[0,1]^n}\abs f$ in the max-row-sum norm.
\item Lemma~\ref{bp:p2:cohomological}: the cohomological equation for vector-valued data in the max norm,
  and the a-priori bound for matrix-valued solutions in the operator norm.
\item Increments on open sets, the first bound of Lemma~\ref{bp:p2:hamiltonian}(e): if $F$ is holomorphic
  on an open $U\subset\C^m$, $\abs{DF}\le A$ on $U$ and $[z,z+w]\subset U$, then
  $\abs{F(z+w)-F(z)}\le A\abs w$ (applied to $F=DX_h$; Lemma~\ref{bp:p1:increments} covers only strips).
\item Realness of the inverse of a real invertible matrix (Lemma~\ref{bp:p2:periodicity}); a one-line
  addition to Lemma~\ref{bp:p1:holomorphic-inverse}.
\item Uniform limits: if holomorphic $f_s$ converge pointwise on an open set and $Df_s$ converge uniformly
  there, the limit is holomorphic with derivative $\lim Df_s$ (Proposition~\ref{bp:p2:limit}; Mathlib
  \Lean{hasFDerivAt\_of\_tendstoUniformlyOn}).
\item Lemma~\ref{bp:p2:error-decay}(a), which strengthens Lemma~\ref{bp:p1:iteration} to a bound
  proportional to $y_0$ and implies it; the two should be unified in this form, as Part~\ref{bp:p1}
  anticipates.
\end{enumerate}

\subsection{Requests to the orchestrator}\label{bp:p2:sec-requests}

\begin{enumerate}[label=(Q\arabic*)]
\item \emph{Lemma~\ref{bp:p1:cohomological}.} State the class in which the solution is unique: all
  functions holomorphic and $1$-periodic on $T_{\rho-\delta}$ with zero average, with no boundedness
  assumed. Lemma~\ref{bp:p2:cohomological}(b) identifies $uv$, and through it $\Omega_K$, with the solution
  by this uniqueness. Also state that for vector-valued $g$ the bound holds in the max norm with no extra
  factor (it does, componentwise).
\item \emph{Lemma~\ref{bp:p1:iteration}.} Replace its conclusion by the form of
  Lemma~\ref{bp:p2:error-decay}(a): the closeness bound of G1 needs $\sum_m y_m/\delta_m^j\lesssim y_0$,
  which the current form ($y_m\le Q^{-1}2^{-p(m+1)}2^{-2^m}$, independent of $y_0$) does not give.
\item \emph{Theorem~\ref{bp:p0:g1} (transcription).} Say that $\abs{D^2X_h}$ is the operator norm of the
  bilinear map (Lean \Lean{iteratedFDeriv}), not HL19's sum of second partial derivatives, and that
  $\abs{Dh}$ is the operator norm, i.e.\ the $\ell^1$ norm of the gradient. Neither affects the proof.
\item \emph{Ledger D-G1-1.} Add: by the Diophantine condition at $k=e_1$ and Lemma~\ref{bp:p2:frequency},
  $\gamma\le\abs\omega\le2n\sigma_K\sigma_B(M+e)$; the bound $\bar\gamma$ is needed only because the
  smallness condition alone does not bound $e$ when $\gamma$ is large.
\item \emph{Ledger:} record \ref{bp:p2:L1}--\ref{bp:p2:L5} as CANDIDATE errata of HL19 and
  \ref{bp:p2:V1}--\ref{bp:p2:V9} as deviations (or fold them into D-G1-1, D-NORM-1, D-COORD-1).
\item \emph{Parts~\ref{bp:p3} and~\ref{bp:p4}.} Lemma~\ref{bp:p2:perturbation} is the stability statement
  STAB of \texttt{docs/PLAN.md} (H2--H4 under $C^0$-small holomorphic perturbations, with explicit strip
  losses); Lemmas~\ref{bp:p2:frame-identities}, \ref{bp:p2:exactness} and~\ref{bp:p2:geometry} hold for
  $E=0$ and give the exact geometry needed in Part~\ref{bp:p3}.
\item \emph{The frozen statement.} No defect found: Theorem~\ref{bp:p0:g1} is proved as stated
  (Theorem~\ref{bp:p2:g1}).
\end{enumerate}
